\documentclass[11pt]{article}
\ifdefined\pdfinfoomitdate\pdfinfoomitdate=1\fi
\ifdefined\pdftrailerid\pdftrailerid{}\fi
\usepackage[T1]{fontenc}
\usepackage[utf8]{inputenc}
\usepackage{lmodern,amsmath,amssymb,amsthm,mathtools,mathrsfs}
\usepackage[margin=26mm]{geometry}
\usepackage{booktabs,array,enumitem}
\usepackage[hidelinks]{hyperref}
\usepackage{microtype,needspace}
\setlength{\emergencystretch}{3em}
\raggedbottom
\setcounter{tocdepth}{1}
\setlist[enumerate]{beginpenalty=10000}
\numberwithin{equation}{section}
\newtheorem{theorem}{Theorem}[section]
\newtheorem{lemma}[theorem]{Lemma}
\newtheorem{proposition}[theorem]{Proposition}
\newtheorem{corollary}[theorem]{Corollary}
\theoremstyle{definition}
\newtheorem{definition}[theorem]{Definition}
\newtheorem{remark}[theorem]{Remark}
\newcommand{\C}{\mathbb C}
\newcommand{\R}{\mathbb R}
\newcommand{\Q}{\mathbb Q}
\newcommand{\Z}{\mathbb Z}
\newcommand{\T}{\mathbb T}
\newcommand{\ii}{\mathrm i}
\newcommand{\dd}{\,\mathrm d}
\newcommand{\eps}{\varepsilon}
\renewcommand{\DH}{\mathcal D_H}
\newcommand{\DL}{\mathcal D_L}
\newcommand{\XH}{\mathcal X_H}
\newcommand{\XL}{\mathcal X_L}
\newcommand{\cO}{\mathcal O}
\newcommand{\norm}[1]{\left\lVert#1\right\rVert}
\DeclareMathOperator{\dist}{dist}
\DeclareMathOperator{\arsinh}{arsinh}
\hypersetup{pdftitle={Boundary jets of the anisotropic square-lattice Ising susceptibility},pdfauthor={Dehao Lin},bookmarksnumbered=true}
\title{Boundary jets of the anisotropic\\square-lattice Ising susceptibility}
\author{Dehao Lin\\
  \small School of Physics, Sun Yat-sen University\\
  \small\href{https://orcid.org/0009-0001-4551-8490}{ORCID 0009-0001-4551-8490}}
\date{21 November 2026}
\makeatletter
\renewcommand{\@maketitle}{\begin{center}
  \vspace*{.5\baselineskip}{\LARGE\@title\par}\vspace{.6\baselineskip}
  {\large\begin{tabular}{c}\@author\end{tabular}\par}
  \vspace{.45\baselineskip}{\normalsize\@date\par}
\end{center}}
\makeatother
\begin{document}
\maketitle
\begin{abstract}
The full high- and low-temperature susceptibilities of the zero-field square-lattice
Ising model are holomorphic on distinct domains in the two complex coupling
variables. Neither admits continuation through any finite point of its
boundary. On each temperature line $K_y=\lambda K_x$ with $\lambda>0$
and $\lambda\ne1$, transverse edge singularities give a natural boundary.
At isotropy the boundary is the unit circle in $s=\sinh(2K)$.
The susceptibility in each phase is differentially transcendental along
every positive-ratio line.
To determine which singularities of the individual particle
integrals survive in the full susceptibility, we bound the derivatives
of their higher-order tail. At isotropy the tail starting at particle
number $n\ge2$ has continuous boundary derivatives on regular arcs whenever
$2j<n^2-3$, where $j$ is the derivative order. At the highest such order
the modulus is $O(h^{1/2})$ for even $n$ and $O(h\log(e/h))$ for odd $n$.
Grouping neighboring lattice directions gives the sharp tail bound. Two families of first resonances indexed
by primes attain the respective moduli. At selected
negative noncollinear edges each odd-particle integral is regular, yet
the one-sided Taylor coefficients $a_j$ of their sum in the normalized
edge coordinate satisfy $\limsup_{j\to\infty}(\log j/j)|a_j|^{1/j}=4/(e\pi^2)$.
\end{abstract}
\begingroup\small
\tableofcontents
\endgroup

\section{Introduction}\label{sec:introduction}

The zero-field susceptibility measures the response of the magnetization
to a uniform magnetic field. Its reduced form is the sum of the connected
two-spin correlation over the square lattice. The high-temperature (HT)
expansion contains the odd particle numbers. Below the critical temperature,
subtracting the squared spontaneous magnetization at each site leaves the
positive even particle numbers in the low-temperature (LT) expansion
\cite{ONGP}.

Guttmann and Enting \cite{GE} found evidence of accumulating singularities
in the anisotropic high-temperature series.
For isotropic coupling $K$, put $s=\sinh(2K)$. Nickel \cite{Nickel}
identified logarithmic branch points of the odd-particle HT integrals
whose union is dense on $|s|=1$. For each fixed even particle number,
Tracy and Widom \cite[Sec.~3, p.~1131]{TW} proved that the integral
has $C^\infty$ boundary values from $|s|>1$ away from its Nickel points.
McCoy and Maillard \cite[Sec.~3.7]{MM} discuss possible cancellation
of these singularities in the full sum. Guttmann, Jensen, Maillard and
Pantone \cite[Sec.~6]{GJMP} used the isotropic HT and LT series to
search for an algebraic differential equation for the full susceptibility.

Theorem~\ref{thm:unified} identifies the boundaries of holomorphy in
the two complex couplings. Fixing $K_y=\lambda K_x$ restricts each
susceptibility to a complex temperature line. For $\lambda>0$,
$\lambda\ne1$, Lemma~\ref{geo:new:slices} shows that the line meets
the selected singularities transversely at a dense subset of its
boundary. At equal coupling $K$ the collinear estimates give the unit
circle in $s=\sinh(2K)$ as the boundary of the HT disk and the LT
exterior. They also determine the last continuous radial derivative
and its modulus of continuity on two families of first resonances.

Suppose the first singular particle integral at a boundary point has
particle number $m$. The higher terms of the same parity begin at $m+2$.
For $2J<(m+2)^2-3$, the derivatives of this tail through order $J$
extend continuously to regular noncollinear and isotropic boundary arcs.
When $\ell$ angles meet, their squared pair differences vanish to
order $\ell(\ell-1)$. The dispersion equation splits a pair factor into two
rational factors. The Cauchy determinant bounds their products, and
zeros at coincident transfer roots give Gaussian decay in the particle number.

At isotropy write $A$ for the larger lattice coordinate in absolute
value. We group neighboring directions on each shell, placing the
centers $\lceil\sqrt A\rceil$ lattice steps apart. Directions in each
interior group share a contour with Gaussian decay in every angle.
Their Fourier sum is a Fej\'er kernel in the total angle.
Its decay away from multiples of $2\pi$, together with the bounds for
the endpoint groups, gives the sharp derivative bound. At a noncollinear endpoint
the Gaussian integral on a plane of fixed total angle records the winding.
The same integral gives the first resonance amplitude.

At a negative noncollinear HT edge the odd momentum constraint keeps
every fixed particle term regular. At the phases of
Theorem~\ref{neg:new:indicator}, the two smallest odd windings
contribute with the same complex phase. Let $a_j$ denote the one-sided
Taylor coefficients of the complete odd-particle sum in the normalized
edge coordinate.
Section~\ref{sec:jet-subsequence} bounds $|a_j|$ at every order and
constructs orders $j_n$ at which one particle number dominates, giving
$|a_{j_n}|^{1/j_n}\sim4j_n/(e\pi^2\log j_n)$.
Lemma~\ref{neg:new:transport} gives the growth rate in physical coordinates.

For each finite $r$ the
susceptibility and its derivatives through order $r$ have continuous
limits on an open boundary arc. An algebraic differential equation
would then give local holomorphic continuation at some point of the
arc, contradicting its natural boundary.

\section{Physical setting and main theorem}
\label{sec:unified}

Let $K_x,K_y$ be independent complex couplings and write
\[
 \begin{gathered}
 z=2K_x=x+\ii y,\qquad w=2K_y=u+\ii v,\\
 S=\sinh z,\quad T=\sinh w,\quad A=\cosh z\cosh w,\quad p=ST.
 \end{gathered}
\]
Define the spectral set by
\[
 \Sigma=\{(K_x,K_y):A=S a_1+T a_2\text{ for some }(a_1,a_2)\in[-1,1]^2\}.
\]
Taking $a_1=\cos(2\pi l/N)$ and $a_2=\cos(2\pi j/N)$ gives the candidate
singularity equation for the anisotropic $N$-particle term in
\cite[(3.23)--(3.24)]{ONGP}. The singular loci of the anisotropic integrals are studied in \cite{BHM}.
In the coupling coordinates $(K_x,K_y)$, set
\begin{align}
 \mathcal D_H={}&\{|y|,|v|<\pi/2,\quad
 \cosh x\tanh|u|<\cos y,\quad \cosh u\tanh|x|<\cos v\},
 \label{eq:domainH}\\
 \mathcal D_L={}&\{x,u>0,\quad
 \cosh x\tanh u>|\cos y|,\quad \cosh u\tanh x>|\cos v|\}.
 \label{eq:domainL}
\end{align}
For real positive noncritical couplings the physical functions are
\begin{align}
 \mathcal X_H&=\sum_{(m,l)\in\Z^2}\langle\sigma_{00}\sigma_{ml}\rangle_H,
 &&0<p<1,\label{eq:defH}\\
 \mathcal X_L&=\sum_{(m,l)\in\Z^2}
       (\langle\sigma_{00}\sigma_{ml}\rangle_+-M^2),
 &&p>1,\qquad M^2=(1-p^{-2})^{1/4}.
 \label{eq:defL}
\end{align}
The fourth root is positive on the real physical region. The origin and
both axes belong to these sums; the low-temperature origin contributes
$1-M^2$. The subtraction is sitewise in the plus state.

A fixed finite positive anisotropy ratio $\lambda$ determines the
temperature line and its physical susceptibility by
\[
 \iota_\lambda(t)=(t/2,\lambda t/2),\qquad
 f_{\nu,\lambda}(t)=\mathcal X_\nu(\iota_\lambda(t)),\qquad
 \chi_{\nu,\lambda}(t)=\frac{t}{2J_x}f_{\nu,\lambda}(t),\quad J_x>0.
\]
Thus $J_y=\lambda J_x$ and the energy-field factor is
$\beta=t/(2J_x)$. Put
\begin{gather*}
 P_\lambda(a)=\cosh a\tanh(\lambda a),\qquad
 Q_\lambda(a)=\cosh(\lambda a)\tanh a,\qquad
 Y_\lambda=\frac{\pi}{2\max(1,\lambda)},\\
 H_\lambda(b)=\min\{P_\lambda^{-1}(\cos b),
               Q_\lambda^{-1}(\cos(\lambda b))\},\quad |b|<Y_\lambda,\\
 L_\lambda(b)=\max\{P_\lambda^{-1}(|\cos b|),
               Q_\lambda^{-1}(|\cos(\lambda b)|)\},\quad b\in\R,\\
 \Omega_{H,\lambda}=\{a+\ii b:|b|<Y_\lambda,\ |a|<H_\lambda(b)\},\qquad
 \Omega_{L,\lambda}=\{a+\ii b:a>L_\lambda(b)\}.
\end{gather*}
The inverses in these formulas are the increasing inverses on $[0,\infty)$.

A function $f$ on a domain $D$ extends through $q\in\partial D$ if a
holomorphic function on some neighborhood $U$ of $q$ agrees with $f$
on a connected component of $D\cap U$ whose closure contains $q$.

On a regular face, $S,T\ne0$, $B=T/S$ is nonreal and $p^2\ne1$.
After exchanging axes if necessary, the spectral representation there
is $A=Sb+T\sigma$, where $-1<b<1$ and $\sigma=\pm1$ is the unique
endpoint coordinate. The edge coordinate $h=A/S-\sigma B$ equals $b$
on the face. A transverse curve with $h'\ne0$ is locally a graph $B=B(h)$.
The selected spectral angles in (iii) have $b=\cos(2\pi r/q)$ with
$q$ an odd prime and $0<r<q/2$. For uniform estimates the graphs are
bounded and holomorphic on a common larger parameter polydisk, with
$|\Im B|\ge\beta>0$ along the edge. The negative HT assertion in (vii)
also holds for any odd denominator $q\ge3$ with $(r,q)=1$.

\begingroup\clubpenalties 6 10000 10000 10000 10000 10000 0\relax
\begin{theorem}\label{thm:unified}
The full-lattice sums \eqref{eq:defH}--\eqref{eq:defL} satisfy the following eight assertions.
\begin{enumerate}[label=\textup{(\roman*)},leftmargin=*]
\item The domains \eqref{eq:domainH}--\eqref{eq:domainL} are exactly the
components of $\C^2\setminus\Sigma$ containing the positive real
high- and low-temperature regions, respectively. They are contractible
and pseudoconvex, with $|p|<1$ on $\mathcal D_H$ and $|p|>1$ on
$\mathcal D_L$.
\item The physical germs define single-valued jointly holomorphic
$\mathcal X_H\in\mathcal O(\mathcal D_H)$ and
$\mathcal X_L\in\mathcal O(\mathcal D_L)$. Their complete odd and
nonzero-even full-site form-factor series converge normally on compact
subsets after every fixed mixed complex derivative.
\item Neither function extends through any finite point of its boundary.
This also excludes a bypass in an unramified Riemann domain reached
by analytic continuation from the physical germ. On a dense subset of
every regular face, nonextension holds on every holomorphic transverse
curve with nonramified edge coordinate. For each fixed selected phase,
the estimates are locally uniform on compact admissible families of
bounded holomorphic curves on common parameter polydisks with $|\Im B|$
uniformly bounded away from zero.
\item For every finite $\lambda>0$,
$\iota_\lambda^{-1}(\mathcal D_\nu)=\Omega_{\nu,\lambda}$.
For each $\nu$, both $f_{\nu,\lambda}$ and $\chi_{\nu,\lambda}$ are
holomorphic on $\Omega_{\nu,\lambda}$ and nonextendable through every
boundary point.
\item At $\lambda=1$, $s=\sinh t$ gives distinct normalized physical germs
\[
 X_H(s)=\frac{(1-s^4)^{1/4}}{s}\sum_{N\ge1,\ N\ {\rm odd}}C_N(1/s),
 \qquad
 X_L(s)=(1-s^{-4})^{1/4}\sum_{N\ge2,\ N\ {\rm even}}C_N(s).
\]
Here $C_N$ is the complete coefficient defined in
\eqref{eq:iso-coefficient} below. The first function is single-valued
on $|s|<1$, with $X_H(0)=1$; the second is single-valued on $|s|>1$,
with onsite term $1-(1-s^{-4})^{1/4}$. Both have the whole unit circle
as a natural boundary. The map $\sinh$ is a biholomorphism from
$\Omega_{H,1}$ to the disk and the universal covering from
$\Omega_{L,1}$ to the exterior, with deck translations
$t\mapsto t+2\pi\ii k$. The physical factor $\beta=t/(2J_x)$ is
interpreted on this cover.
\item Let $m\ge2$ be an integer cutoff. The complete tail over
$N\ge m+2$, $N\equiv m\pmod2$, satisfies the collar estimates
through derivative order $J$ whenever
\begin{equation}
                 2J<(m+2)^2-3.
                 \label{eq:common-budget}
\end{equation}
The derivatives are continuous and uniformly bounded on the closed
regular noncollinear collars of Theorem~\ref{edge:new:budget}.
At isotropy, with $E=s+s^{-1}$, they extend continuously to each compact
subinterval of $(-2,0)\cup(0,2)$ from the physical side, with bounds
summable in $N$. Write $n=m+2$ for the first tail index and let $j_*$
be the greatest integer with $2j_*<n^2-3$. On shorter closed isotropic
collars, the $j_*$th derivative of this tail has modulus $O(h^{1/2})$
for even $n$ and $O(h\log(e/h))$ for odd $n$, where $h$ is distance
in $E$. Its lower derivatives are Lipschitz.
For each physical slice and each finite order $r$, the complete
physical function has continuous finite jets through order $r$ from a
one-sided collar of a nonempty relatively open real-analytic boundary
arc $I_r$. Negative odd noncollinear edges have continuous jets of
every order on their entire regular arc.
\item On the distinct-angle prime families of
Lemma~\ref{iso:new:radial:arithmetic}, let $n$ be the first resonant
particle number and $k=\lfloor(n^2-1)/2\rfloor$. In the radial
parameter $\epsilon\downarrow0$ of Theorem~\ref{iso:new:radial:full},
the first nonanalytic term is proportional to $\epsilon^{k-1/2}$
in LT and to $\epsilon^{k-1}\log\epsilon$ in HT. The theorem gives
the nonzero amplitudes in the physical angular measure, the constants
of integration and the remainders for the full susceptibility.
The last continuous radial derivative is sharply H\"older $1/2$
in LT and sharply log-Lipschitz in HT.

At each selected noncollinear negative odd edge write $B=T/S$ and
$h=A/S+B$ in the transfer chart, and let $b$ be the value of $h$ at
the boundary point. In the edge coordinate $\tau=(h-b)/B$ the complete
boundary jet has indicator $4/(e\pi^2)$, with an upper bound at every
order and a matching lower subsequence for each family member.
Lemma~\ref{neg:new:transport} gives the exact transport under physical
factors and nonramified coordinates.
\item Separately for each phase, $X_H(s)$ and $X_L(s)$ are
differentially transcendental over $\C(s)$ and their corresponding
physical germs under $s=\sinh t$ are differentially transcendental
over $\C(t)$. For every fixed finite $\lambda>0$, including irrational
$\lambda$, both $f_{\nu,\lambda}$ and $\chi_{\nu,\lambda}$ are
differentially transcendental over $\C(t)$, for each
$\nu\in\{H,L\}$. Thus no nonzero polynomial differential identity
of any finite order or degree holds.
\end{enumerate}
\end{theorem}

\endgroup

\paragraph{Proof outline.}
The angular integral in Section~\ref{sec:physical} represents the
physical susceptibility throughout each domain. The tail estimates of
Sections~\ref{edge:new:section} and~\ref{iso:new:section} isolate the
singular part of the first resonant integral at positive and even
edges. At noncollinear negative odd edges the obstruction comes from the growth
of the Taylor coefficients of the full sum
(Section~\ref{neg:new:section}). The same tail estimates give the
continuous boundary derivatives used in the proof of differential
transcendence in Section~\ref{da:new:section}.


\section{Geometry of the physical continuations}
\label{geo:new:section}

Put $z=2K_x=x+iy$, $w=2K_y=u+iv$, and
\[
 S=\sinh z,\quad T=\sinh w,\quad A=\cosh z\cosh w,\quad p=ST.
\]
In these coordinates the spectral set is
\[
 \Sigma=\{(z,w):A=Sa+Tb\text{ for some }(a,b)\in[-1,1]^2\}.
\]
The representing square is compact, so $\Sigma$ is closed. Its
frontier is contained in the set of incidences for which one of the
two coefficients is an endpoint of $[-1,1]$.

\begin{lemma}[Endpoint incidence determines the spectral frontier]
\label{geo:new:endpoint}
Define
\[
 F_1=\cosh^2u\cos^2y-\cosh^2x\sinh^2u,\qquad
 F_2=\cosh^2x\cos^2v-\cosh^2u\sinh^2x.
\]
The set $E=\{F_1F_2=0\}$ is exactly the set admitting a
spectral representation with $a$ or $b$ an endpoint of $[-1,1]$.
In particular $\partial\Sigma\subset E\subset\Sigma$.
\end{lemma}
\begin{proof}
For $S\ne0$ the quantities $\tanh(z/2)$ and its reciprocal
are finite.  With $\epsilon\in\{1,-1\}$,
\[
 r_\epsilon=e^{\epsilon w}\tanh(z/2),\qquad
 \frac{r_\epsilon+r_\epsilon^{-1}}2=\frac{A-\epsilon T}{S},
 \qquad
 |\tanh(z/2)|^2=\frac{\cosh x-\cos y}{\cosh x+\cos y}.
\]
The numerator and denominator of the last ratio are positive.
Under the Joukowski map $r\mapsto(r+r^{-1})/2$, the preimage
of $[-1,1]$ is the unit circle: the roots of $r^2-2ar+1=0$ for
$a\in[-1,1]$ have unit modulus.  Therefore a representation
$A=Sa+\epsilon T$, $a\in[-1,1]$, is equivalent to
\[
 \cos y=\epsilon\cosh x\tanh u.
\]
The union over $\epsilon$ is $F_1=0$.  If $S=0$ then
$x=0$, $y\in\pi\mathbb Z$, $F_1=1$, and $A^2-T^2=1$,
so neither condition occurs.  Exchanging the variables handles
$F_2$ and the endpoint $a$.

It remains to relate these incidences to the frontier. At a point
of $\partial\Sigma$, consider the real-linear map
$(a,b)\mapsto Sa+Tb$.  If it has rank two, a representing pair
strictly inside the square would persist there under parameter
perturbation, making the point interior to $\Sigma$.  Hence
one coordinate is an endpoint.  If its rank is smaller, move a
representing pair along a nonzero kernel direction to the boundary
of the square.  This again produces an endpoint representation.
\end{proof}

\begin{proposition}[Exact physical domains and branches]
\label{geo:new:domains}
The complementary components containing the positive real HT and
LT regions are respectively
\begin{align*}
 \mathcal D_H={}&\{|y|,|v|<\pi/2,\quad
   \cosh x\tanh|u|<\cos y,\quad
   \cosh u\tanh|x|<\cos v\},\\
 \mathcal D_L={}&\{x,u>0,\quad
   \cosh x\tanh u>|\cos y|,\quad
   \cosh u\tanh x>|\cos v|\}.
\end{align*}
They are contractible and pseudoconvex, and satisfy $|p|<1$
and $|p|>1$, respectively.  The binomial branches of
$(1-p^2)^{1/4}$ and $(1-p^{-2})^{1/4}$ are thus single valued
on their respective domains, with the positive physical values.
On $\mathcal D_H$ the only zeros of $S$ and $T$ are $z=0$
and $w=0$; neither vanishes on $\mathcal D_L$.
\end{proposition}
\begin{proof}
Write $r_1(a,b)=\cosh a\tanh b$ and
$r_2(a,b)=\cosh b\tanh a$ for $a,b\geq0$.  Both are
nondecreasing in either variable, and
\[
 r_1^2-1=\frac{\sinh^2a\sinh^2b-1}{\cosh^2b},\qquad
 r_2^2-1=\frac{\sinh^2a\sinh^2b-1}{\cosh^2a}.
\]
\emph{The high-temperature component.}
Dilating all four real coordinates toward zero preserves
$\mathcal D_H$. Its closure is given by the corresponding non-strict
inequalities. Dilation by any positive factor smaller than one moves
every point satisfying these inequalities into $\mathcal D_H$. For an active first
inequality this is immediate if $u\ne0$; if $u=0$, activity
forces $|y|=\pi/2$, and the cosine strictly increases under
dilation.  The second inequality has the exchanged argument.
Thus $\mathcal D_H$ is star shaped, and its finite boundary is
in $E$.

\emph{The low-temperature component.}
Increasing both real parts makes all the non-strict inequalities
strict, including $x,u\geq0$. Consequently the closure is given
by these non-strict inequalities, and its finite
boundary is also in $E$.  Fix $R>0$ with $\sinh R>1$.
Increase $x,u$ to $\max(x,R),\max(u,R)$, move them to $R,R$
within $[R,\infty)^2$, and then move $y,v$ to zero.
Throughout the last two motions both $r_i\geq\sinh R>1$.
This is a contraction of $\mathcal D_L$.

\emph{Component identification and branches.}
Each displayed domain is connected and avoids $E$. Since
$\partial\Sigma\subset E$, its intersection with $\Sigma$
is both relatively open and closed.  Zero lies outside $\Sigma$,
and so does the real LT point $(R,R)$: there
\[
 (A-S-T)(A+S+T)=(p-1)^2>0,
\]
with $A+S+T>0$.  The two domains are therefore disjoint from
$\Sigma$.  Their boundaries lie in $\Sigma$, so each is an
entire complementary component.  On the positive real quadrant
the preceding formulas for $r_i^2-1$ identify the HT region
$0<p<1$ and LT region $p>1$.

The identity
\[
 |p|^2=(\cosh^2x-\cos^2y)(\cosh^2u-\cos^2v)
\]
gives the modulus statements directly: for $F_1,F_2>0$ the
two factors are smaller than $\cosh^2x/\cosh^2u$ and its
reciprocal, while for $F_1,F_2<0$ both inequalities reverse.
The statements about branches and zeros follow.

\emph{Pseudoconvexity.}
Put $m(K)=\min_{a,b\in[-1,1]}|A-Sa-Tb|$.
On every complementary component it is continuous and positive,
and $-\log m$ is the continuous locally bounded supremum of the
pluriharmonic functions $-\log|A-Sa-Tb|$.  Applying the submean
inequality on each complex line proves that $-\log m$ is
plurisubharmonic.  Thus
\[
 |K_x|^2+|K_y|^2+\max\{0,-\log m(K)\}
\]
is a continuous plurisubharmonic exhaustion.  It diverges at
infinity and at every finite boundary point.  A bounded sublevel
has $m$ bounded away from zero, so its limit points remain in
the same component.  This verifies pseudoconvexity.
\end{proof}

\begin{corollary}[The whole spectral set]\label{geo:whole}
One has $\Sigma=\{F_1F_2\le0\}$ and $\partial\Sigma=\{F_1F_2=0\}$.
The components with $F_1,F_2>0$
are the translates of $\mathcal D_H$ by $(\ii m\pi,\ii k\pi)$ in
$(z,w)$, $m,k\in\Z$. Those with $F_1,F_2<0$ are the four images of
$\mathcal D_L$ under independent sign changes of $z,w$.
\end{corollary}
\begin{proof}
The indicated translations and reflections preserve $\Sigma$ by
sign changes in its representing square. The strict equal-sign
inequalities give exactly the listed domains by
Proposition~\ref{geo:new:domains}. If $F_1>0>F_2$, hold $x,u$
fixed and move $y$ to its nearest $m\pi$ and $v$ to its nearest
$(k+1/2)\pi$. Along these segments $|\cos y|$ increases and
$|\cos v|$ decreases, preserving the signs. At the endpoint
$A=T(-1)^m\cosh x\tanh u$, whose displayed coefficient has
absolute value less than one. This is a spectral representation
with the coefficient of $S$ equal to zero. Since no point on the
path can meet $\partial\Sigma$ by Lemma~\ref{geo:new:endpoint},
the starting point is spectral. Exchange the axes for $F_2>0>F_1$;
the zero set already belongs to $\Sigma$ by the same lemma.

Every zero is approachable from an equal-sign region. At a double
zero, increasing both $|x|,|u|$ makes both signs negative. At
$F_1=0>F_2$ the same move works. At $F_1=0<F_2$ decrease $|u|$
if $u\ne0$; when $u=0$ increase $|\cos y|$ from zero instead.
The other cases follow by exchanging axes.
\end{proof}

\begin{lemma}[Canonical transfer root]
\label{geo:new:root}
Off $\Sigma$, for each real $q$, the equation
\[
 S\rho^2-2(A-T\cos q)\rho+S=0
\]
has a unique root of modulus less than one, interpreted as zero
when $S=0$.  This root is holomorphic in the couplings and locally
jointly holomorphic in a complex angle.  On any compact parameter
set off $\Sigma$ it has a uniform modulus gap below one on the
whole real angular circle.  In both positive real physical phases
it is the positive decaying root.
\end{lemma}
\begin{proof}
For $S\ne0$, $W=(A-T\cos q)/S$ lies outside $[-1,1]$.
The two roots are reciprocal, distinct, and neither has modulus
one.  Exactly one is inside; its implicit-function branches agree
on overlaps by uniqueness.  When $S=0$, $A-T\cos q\ne0$,
so the implicit-function theorem extends the same root through
$\rho=0$.  Compactness gives the modulus gap and a common thin
complex angular neighborhood.  At positive noncritical couplings
the identity above gives $A-S-T>0$, so $W>1$ and $0<\rho<1$.
\end{proof}

At a regular face one spectral coefficient is an endpoint and the
other is a local coordinate along the spectral interval. Such faces
are dense in the boundary of each physical component.

\begin{proposition}[Regular faces, coordinates, and dense phases]
\label{geo:new:faces}
Points of either physical boundary with $x,u\ne0$ and exactly
one zero among $F_1,F_2$ are dense in that boundary.  Each is
a regular edge: $S,T\ne0$, $\Im(\overline S T)\ne0$, its
unique representing pair has exactly one endpoint, the other
coordinate is in $(-1,1)$, and $p^2\ne1$.

Where the endpoint is $b=\epsilon$, the variables
\[
 C=A/S,\qquad B=T/S,\qquad h=C-\epsilon B
\]
give a holomorphic chart in which the boundary is $\Im h=0$,
$\Re h\in(-1,1)$, and the physical domain is its entire local
complementary side.  When the endpoint is $a=\epsilon$, exchange
the transfer axes.  On every regular face, points with free
coefficient $\cos(2\pi k/q)$, $q$ an odd prime and
$1\leq k\leq(q-1)/2$, are dense.
\end{proposition}
\begin{proof}
At $F_1=0$ use Lemma~\ref{geo:new:endpoint} to write
$A=Sa+\epsilon T$, $a\in[-1,1]$.  Since $x,u\ne0$, both $S,T$ are nonzero. By
Lemma~\ref{geo:new:endpoint} the sign $\epsilon$ is unique, and
$F_2\ne0$ gives $-1<a<1$. If $S,T$ were collinear, their real
kernel line through $(a,\epsilon)$ would enter the representing
square and leave at another endpoint. This would give either
$F_2=0$ or the other sign of $\epsilon$. Also $p=\pm1$ would give
a corner representation, by $A^2-(S+\delta T)^2=(p-\delta)^2$
with $\delta=p$. These contradictions give all the strict conditions.
Exchange the axes for $F_2=0$.

\emph{Density of regular faces.}
Fix real-part magnitudes $a,b$ with
$\sinh a\sinh b<1$.  The HT closure has angular fiber
\[
 |y|\leq\arccos r_1(a,b),\qquad
 |v|\leq\arccos r_2(a,b).
\]
At product one both intervals collapse to zero.  Any HT boundary
point has an active face.  Approximate its magnitudes by positive
$a_j,b_j$ with product below one, keep an active angular coordinate
at its endpoint, and place the other strictly inside its interval.
Choose signs to approach the prescribed coordinates.  This gives
exactly one zero $F_i$ and nonzero real parts, also at strip tips,
intersections, and collapsed fibers.

At an LT boundary point $\sinh x\sinh u\leq1$; otherwise both
inequalities are strict.  At least one cosine inequality is active,
also when a real part is zero.  Approximate $x,u$ by positive
$x_j,u_j$ with product below one.  If the first face is active,
choose $|\cos y_j|=r_1(x_j,u_j)$ tending to the prescribed value
and choose $|\cos v_j|<r_2(x_j,u_j)$ tending to the other value.
The local cosine intervals permit this also at their zero or
unit endpoints.  Exchange the roles when the second is active.
This proves density at every exceptional stratum as well.

\emph{Edge coordinates and resonant phases.}
Direct differentiation gives
\[
 \det\frac{\partial(C,B)}{\partial(z,w)}=\frac{p^2-1}{S^3},
 \qquad \Im(\overline S T)=|S|^2\Im B.
\]
In this chart the spectral set is the parallelogram
$C\in[-1,1]+B[-1,1]$.  Near the interior of its endpoint
edge, the representing coefficient satisfies
\[
 b-\epsilon=\frac{\Im(C-\epsilon B)}{\Im B},
\]
while the other coefficient stays strictly interior.  One local
side is spectral and the other is connected and complementary;
the latter meets, and therefore lies in, the physical component.
For fixed $B$ the variable $h$ ranges over an open real interval
on the face.  The mesh of the prime-angle grids tends to zero,
proving the last assertion.
\end{proof}

\subsection{Physical temperature lines at arbitrary positive anisotropy}

To determine the physical domains on each prescribed temperature
line, define for $\lambda>0$ the odd real-analytic increasing
bijections of $\mathbb R$
\[
 P_\lambda(a)=\cosh a\tanh(\lambda a),\qquad
 Q_\lambda(a)=\cosh(\lambda a)\tanh a.
\]
Their derivatives are strictly positive; for example
$P_\lambda'(a)=\sinh a\tanh(\lambda a)+
\lambda\cosh a\operatorname{sech}^2(\lambda a)>0$.
Let $z_c>0$ solve $\sinh z_c\sinh(\lambda z_c)=1$ and put
\[
 Y_\lambda=\frac\pi{2\max(1,\lambda)},\quad
 H_\lambda(b)=\min\{P_\lambda^{-1}(\cos b),
                         Q_\lambda^{-1}(\cos(\lambda b))\},
 \quad |b|<Y_\lambda,
\]
\[
 L_\lambda(b)=\max\{P_\lambda^{-1}(|\cos b|),
                         Q_\lambda^{-1}(|\cos(\lambda b)|)\}.
\]
Substitution of $(z,w)=(t,\lambda t)$, $t=a+ib$, in
Proposition~\ref{geo:new:domains} gives exactly
\begin{equation}\label{geo:new:pullbacks}
 \Omega_{H,\lambda}=\{|b|<Y_\lambda,\ |a|<H_\lambda(b)\},
 \qquad
 \Omega_{L,\lambda}=\{a>L_\lambda(b)\}.
\end{equation}
Here $P_\lambda(z_c)=Q_\lambda(z_c)=1$, so
$0\leq L_\lambda\leq z_c$.  The HT profile is positive
on its interval and tends to zero at both tips.  Its domain is
star shaped by dilation; the LT domain is homeomorphic to a
half-plane under $(a,b)\mapsto(a-L_\lambda(b),b)$.
Both are therefore connected and simply connected.

\begin{lemma}[Regular points on each prescribed line]
\label{geo:new:slices}
For every fixed $\lambda>0$, $\lambda\ne1$, all but a locally
finite subset of each boundary in \eqref{geo:new:pullbacks}
are regular joint edges.  The edge coordinate has nonzero
derivative on the prescribed temperature line at every such point.
The selected prime phases are consequently dense in every
regular boundary arc and in the whole slice boundary.
\end{lemma}
\begin{proof}
Define the signed profiles by
\[
 R(b)=P_\lambda^{-1}(\cos b),\quad
 U(b)=Q_\lambda^{-1}(\cos(\lambda b)),\quad
 E_\lambda(b)=R(b)^2-U(b)^2.
\]
The signed profiles are real analytic on the whole real line.
At $b=Y_\lambda=\pi/(2\max(1,\lambda))$, exactly one of
$\cos b$ and $\cos(\lambda b)$ is zero and the other is positive.
As $P_\lambda(0)=Q_\lambda(0)=0$, $E_\lambda(Y_\lambda)\ne0$.
Its zeros are therefore locally finite.

Away from those contacts and the two HT tips, each boundary point
has nonzero real part and exactly one active inequality.
Proposition~\ref{geo:new:faces} supplies its regular edge chart.
On an $F_1$ face write
$r=e^{\epsilon\lambda t}\tanh(t/2)$, so $|r|=1$,
$h=(r+r^{-1})/2$ and $(\log r)'=\epsilon\lambda+1/\sinh t$.
For $t=a+ib$, the endpoint identity gives
\[
 \epsilon\Re(1/\sinh t)
  =\frac{\sinh a\cosh a\tanh(\lambda a)}{|\sinh t|^2}\ge0.
\]
Hence $|h'(t)|\ge\lambda\sqrt{1-h^2}$. Exchanging axes gives
$|h'(t)|\ge\sqrt{1-h^2}$ on an $F_2$ face, using
$r=e^{\epsilon t}\tanh(\lambda t/2)$. The free coefficient
therefore runs through an open interval on that line. Prime
phases are dense there, and deleting the locally finite contacts
and the tips leaves a dense subset of the whole slice boundary.
\end{proof}

On a regular HT $F_1$ face the endpoint is
$\epsilon=\operatorname{sign}u$; on an LT $F_1$ face it is
$\operatorname{sign}(\cos y)$.  Exchange coordinates on $F_2$.
Thus the left HT boundary of every unequal positive-ratio slice
has negative endpoint, and its right boundary has positive
endpoint.  A regular selected point has $t_0\ne0$.  The map
$(t,\lambda)\mapsto(z,w)=(t,\lambda t)$ has determinant $t_0$.
The determinant and the nonzero edge derivative give $B=B(h,\lambda)$
on a product polydisk by the holomorphic implicit-function theorem.
On a smaller polydisk, $|\lambda-1|$ and $|\Im B|$ have positive lower
bounds, the sign of $\Im B$ is fixed, and the physical factors are
nonzero. The larger polydisk supplies common Cauchy bounds for the
parameter family.

\subsection{The spectrum on a fixed-ratio line}\label{sec:slice-spectrum}

Write $\Sigma_\lambda=\{t:(t/2,\lambda t/2)\in\Sigma\}$. Since
$P_\lambda(z_c)=Q_\lambda(z_c)=1$, monotonicity gives
$|\Re t|\le z_c$ on $\Sigma_\lambda$. Near $t=z_c$ at height $b$,
the width between the two spectral edges is
\[
 \tfrac12|\lambda\coth(\lambda z_c)-\coth z_c|\, b^2+O(b^4).
\]
Indeed the edges have real parts $z_c-b^2/(2P'_\lambda(z_c))+O(b^4)$
and $z_c-\lambda^2 b^2/(2Q'_\lambda(z_c))+O(b^4)$.
For $h=\sinh z_c$ and $c=\cosh z_c$ the two derivatives are
$h(1+\lambda h)/c$ and $(1+\lambda h)/(hc)$. Their substitution
gives the displayed coefficient, positive for $\lambda\ne1$ by
strict increase of $v\coth v$ for $v>0$.

On the imaginary axis the spectral condition is
$\cos b\cos(\lambda b)=0$. Simultaneous zeros occur precisely when
$\lambda=r/s$ in lowest terms with both entries odd, and then
$b=(2j+1)\pi s/2$ with $j\in\Z$. The spectrum is invariant under
sign and conjugation, and $\Sigma_{1/\lambda}=\lambda\Sigma_\lambda$
by axis exchange. For rational $\lambda=r/s$, $\ii\pi s$ is a
period by the two incidence tests. For irrational $\lambda$ there
is no nonzero period. The strip bound forces every period to be
imaginary. The two progressions on the imaginary axis have
incommensurable steps, so a translate of one meets the other at
most once. A period would have to preserve both progressions.

For each of the two parities $\Sigma_\lambda$ is the closure of the
zeros of $A-S\cos(2\pi j/N)-T\cos(2\pi k/N)$ with $N=p$ or
$2p$, respectively, for odd primes $p$. For any representing
pair in $[-1,1]^2$, choose grid pairs converging to it. The points
$(j,k)=(0,0)$ and, in even order, $j+k=N/2$ may be omitted.
The corresponding entire functions converge locally uniformly to
$A-Sa_1-Ta_2$, which equals one at $t=0$. Rouch\'e's theorem on
small circles about any zero proves the assertion. Since $S,T$
cannot both vanish at a spectral point, the approximating pairs can
also be chosen off its representing affine equation. Thus
$\Sigma_\lambda$ has no isolated points.

At a regular negative HT endpoint choose the transfer axis so that
$a_2=-1$ and take $a_1=b_0=\cos(2\pi r/q)$, where $q$ is an odd
prime. Throughout the HT lens $\Re(T\overline S)>0$ except at zero,
as follows by expanding its two nonnegative terms
$\sinh a\,\sinh(\lambda a)\cos b\cos(\lambda b)$ and
$\cosh a\,\cosh(\lambda a)\sin b\sin(\lambda b)$. Thus $\sqrt B$
has its principal value, also on the regular boundary by continuity.
Set $D(t;a_1,a_2)=A-Sa_1-Ta_2$. At a selected endpoint $t_0$,
$D_t(t_0;b_0,-1)\ne0$ by Lemma~\ref{geo:new:slices}.
For odd multiples $n$ of $q$, $-\cos(\pi/n)$ is an adjacent
$n$-grid value. The implicit function theorem gives the displaced root
\[
 t_n=t_0+\frac{T(t_0)}{D_t(t_0;b_0,-1)}
                  \frac{\pi^2}{2n^2}+O(n^{-4}).
\]
For large $n$ the two directions remain independent and the representing
pair is strictly interior, so $t_n$ is an interior point of the spectrum.

When $\Re t\ne0$, both $S$ and $T$ are nonzero. If they are collinear
at a spectral point, the representing segment cannot meet opposite
edges of the same coordinate: the endpoint identities would give
$\Re t=0$ or $\lambda\Re t=0$. Hence it meets edges of both
coordinates, including the case of a single corner. Thus $F_1=F_2=0$
and the product formula for $|p|$ gives $|p|=1$. Conversely, if
$T/S=r\in\R$ and $|p|=1$, write $S=|S|e^{\ii\theta}$. Then
$(A/S)^2=1+r^2+2|r|\cos(2\theta)$, so $A/S$ is real of modulus
at most $1+|r|$ and a representing pair exists. At unequal ratios
these points lie among the locally finite profile intersections of
Lemma~\ref{geo:new:slices}.

\subsection{Exceptional points and the isotropic covering}

\begin{lemma}[Inward deformation and componentwise continuation]
\label{geo:new:collar}
For either physical joint domain or either slice domain, if
$P\in\partial D$ and $U$ is a neighborhood of $P$, there is
a neighborhood $V\subset U$ of $P$ such that $D\cap V$ lies
in a single component of $D\cap U$.  Every component of
$D\cap U$ accumulating at $P$ contains $D\cap V$.
Consequently a dense set of local obstructions on the full domain side
excludes componentwise extension at every finite frontier point.
\end{lemma}
\begin{proof}
For HT use $H_s(Q)=(1-s)Q$.  For joint LT increase both real
parts by $s$; on an LT slice use $t\mapsto t+s$.
The domain and closure proofs show that these deformations
preserve $D$ and move every boundary point inside for $0<s<1$.
Choose a small fixed $s>0$ with $H_v(P)\in U$ for $0\leq v\leq s$,
and a ball about $H_s(P)$ contained in $D\cap U$.
Continuity uniformly for $v\in[0,s]$ gives a sufficiently small
$V$ for which all paths $H_v(Q)$, $Q\in D\cap V$, remain
in $D\cap U$ and end in that ball.  Join endpoints in the
ball to connect any two points of $D\cap V$. Every accumulating
component meets $D\cap V$, so it is the same component.

If an extension at $P$ agrees on an accumulating component,
it therefore agrees on $D\cap V$.  Choose a dense obstructing
point of $\partial D$ inside $V$ and restrict the proposed
extension to its neighborhood.  This contradicts its defining
obstruction.
\end{proof}

Consider a continuation path in an unramified Riemann domain over
the coupling space. As long as its projection lies in $D$, uniqueness
identifies the germ with the single-valued function on $D$. If the
projection leaves $D$, its compact image meets the finite boundary
at the first exit. A local biholomorphic chart there would extend the
function through that point on an accumulating component, contrary
to the preceding lemma.

\begin{lemma}[Disk and exterior coordinates]
\label{geo:new:cover}
At $\lambda=1$, put $s=\sinh t$.  The exact spectral preimage is
$\{|\sinh t|=1\}$, and
\[
 \Omega_{H,1}=\{|b|<\pi/2,\ |a|<\operatorname{arsinh}(\cos b)\},
 \qquad
 \Omega_{L,1}=\{a>\operatorname{arsinh}(|\cos b|)\}.
\]
The map $\sinh$ is a biholomorphism from $\Omega_{H,1}$ onto
$|s|<1$ and a universal covering from $\Omega_{L,1}$ onto
$|s|>1$, with deck translations $t\mapsto t+2\pi i k$.
\end{lemma}
\begin{proof}
Now $S=T=s$, $A=1+s^2$.  For $s\ne0$ the unsquared
spectral condition is $s+s^{-1}\in[-2,2]$, precisely $|s|=1$;
$s=0$ is nonspectral.  The profiles follow from
$P_1=Q_1=\sinh$.

The identity $|\sinh(a+ib)|^2=\sinh^2a+\sin^2b$
maps the HT lens into the disk and its boundary into the circle.
Equal hyperbolic sines require either
$t_1=t_2+2\pi i k$ or $t_1=-t_2+(2k+1)\pi i$.
In the central strip only the first alternative with $k=0$
is possible.  The derivative has no zero there.  The lens is
bounded, so its image is both open and relatively closed in
the disk, and contains zero.  It is the entire disk.

For the exterior put $\omega=e^t$.  The equation
$\omega-\omega^{-1}=2s$ has exactly one root with $|\omega|>1$
for $|s|>1$: the root moduli have product one, and a unit root
would give $s\in i[-1,1]$.  The outside root varies holomorphically,
using $\sqrt{s^2+1}=s\sqrt{1+s^{-2}}$ with the binomial branch.
Its image is
\[
 \{re^{ib}:r>\exp(\operatorname{arsinh}|\cos b|)\},
\]
by the same modulus identity.  This root map is a biholomorphism.
Its full exponential preimage is the simply connected LT epigraph,
so exponentiation is its universal covering and has the stated
deck group.
\end{proof}

The reduced LT susceptibility descends once its termwise invariance
under these deck translations has been proved in
Section~\ref{sec:physical}. After descent, an extension of either
descended function through a unit-circle point would compose with the
entire map $\sinh$ to extend the physical slice. This implication also
holds at the ramified boundary values $s=\pm i$. The physical factor
$\beta=t/(2J_x)$
is single valued on the physical temperature domain and nonzero
on its boundary: zero is HT interior and is outside the LT
closure.  It therefore preserves boundary extendability there;
in the LT exterior it remains a function on the covering.

\section{The full-lattice angular integral}\label{sec:physical}

At positive noncritical couplings, use the form-factor expansion for
the two-spin correlation as written by Orrick, Nickel, Guttmann and Perk.
Their $s,s',\cot\alpha$ are $S,T,\sqrt{T/S}$, and their modulus is $p$
in HT and $p^{-1}$ in LT \cite[(1.2), (1.6), (2.9)]{ONGP}.
Their $x_i$ is the positive decaying root
$\rho_i=e^{-\gamma_i}$ \cite[(3.3), (3.14)--(3.16)]{ONGP}, where
\[
 B=T/S,\qquad C=A/S,\qquad \cosh\gamma_i=C-B\cos q_i .
\]
Let $Z=\prod_i\rho_i$, $d\lambda_N=\prod_{i=1}^N dq_i/(2\pi)$
and put
\[
 P_{ij}=B\frac{\sin^2((q_i-q_j)/2)}{\sinh^2((\gamma_i+\gamma_j)/2)},
 \qquad
 w_N(q)=\frac{B^{N/2}}{N!}\prod_i\frac1{\sinh\gamma_i}\prod_{i<j}P_{ij}.
\]
Here $P_{ij}$ is the square of the pair factor in
\cite[(3.19)--(3.20)]{ONGP}. The plus-state magnetization $M$ vanishes
in HT and satisfies $M^2=(1-p^{-2})^{1/4}$ in LT \cite[(2.2)]{ONGP}.
Write $e_H=(1-p^2)^{1/4}p^{-1/2}$,
$e_L=(1-p^{-2})^{1/4}$, and let $\mathcal N_H,\mathcal N_L$ be the
positive odd and positive even indices respectively. For the
homogeneous zero-field Gibbs state $\mu$ in the site formula, setting
the particle fugacity to one in \cite[(4.1)--(4.2)]{ONGP} gives
\begin{equation}\label{eq:real-site}
 C(a,b)-M^2=e_\nu\sum_{N\in\mathcal N_\nu}
 \int_{\mathbb T^N} w_N(q)\,Z^{|a|}
                 \cos\!\bigl(b\textstyle\sum_i q_i\bigr)\,d\lambda_N,
 \qquad C(a,b)=\mu(\sigma_0\sigma_{(a,b)}).
\end{equation}
The LT zero-particle term is $M^2$, and $C(-a,-b)=C(a,b)$.
The site $(0,0)$ contributes $1-M^2$. In $w_N$ the one-particle and
pair factors give a total power $N/2+N(N-1)/2=N^2/2$ of $B$.
The finite-cylinder weight in \cite[(39)]{BL} has the same power.
With $t_x=\tanh K_x$ and $t_y=\tanh K_y$, its anisotropy ratio
$t_y(1-t_x^2)/(t_x(1-t_y^2))$ equals $T/S$.

For $q_N=-\sum_{i<N}q_i$ modulo $2\pi$, set
$d\mu_N=\prod_{i<N}dq_i/(2\pi)$ and define
\begin{equation}\label{eq:complete-angular}
 C_N(C,B)=\int_{\sum q_i=0} w_N(q)\frac{1+Z}{1-Z}\,d\mu_N .
\end{equation}
Summing the correlation in \eqref{eq:real-site} over the lattice gives
\begin{align}
 \mathcal X_H&=(1-p^2)^{1/4}p^{-1/2}
                    \sum_{N\ {\rm odd}} C_N(C,B),\label{eq:physical-formsH}\\
 \mathcal X_L&=(1-p^{-2})^{1/4}
                    \sum_{N\ge2,\ N\ {\rm even}} C_N(C,B)
                    \label{eq:physical-forms}
\end{align}
with the normalization of \cite[(1.8)--(1.10), (3.13)]{ONGP}.
The site sums and their Fourier evaluation are justified in
Section~\ref{sec:lattice-sum}.

\subsection{The two factors of an angular pair}

Put $z_i=e^{\ii q_i}$, $u_i=\rho_i$ and $H(v,w)=(v-w)/(1-vw)$.
Subtracting the dispersion equations gives
$(u_i-u_j)(1-u_i u_j)=-B(u_i u_j/z_i z_j)(z_i-z_j)(1-z_i z_j)$.
Consequently
\begin{equation}\label{eq:split-pair}
 P_{ij}=-\frac{Bu_i u_j}{z_i z_j}
        \left(\frac{z_i-z_j}{1-u_i u_j}\right)^2
       =H(z_i,z_j)H(u_i,u_j),
\end{equation}
where the first expression fills removable values in the second.

On each contour used below the complete pair product is estimated by
Lemma~\ref{lem:analytic-det} in Appendix~\ref{app:forms}. In particular,
one may divide out the squared differences on any fixed block of angles.

\subsection{A symmetric site integral}

For $a,b\in\mathbb Z_{\ge0}$ and $D(v,w)=A-S(v+v^{-1})/2-T(w+w^{-1})/2$, put
\begin{equation}\label{eq:double-site}
 J_N(a,b;A,S,T)=\frac1{N!}\int
 \prod_{i<j} H(v_i,v_j)H(w_i,w_j)
 \prod_{i=1}^N \frac{v_i^a w_i^b}{D(v_i,w_i)}
      \frac{\dd v_i}{2\pi\ii v_i}\frac{\dd w_i}{2\pi\ii w_i}.
\end{equation}
All circles have the same radius $r<1$. Off $\Sigma$, $D$ has a
uniform gap on the unit tori on each compact parameter set. It keeps
this gap when all radii lie between $r$ and $1$ for $r$ sufficiently
close to one. The integral is independent of that choice of $r$
by Cauchy's theorem with the circles strictly inside the unit disk.
Swapping the two sets gives $J_N(a,b;A,S,T)=J_N(b,a;A,T,S)$.
For $S=T=1$ and even $N$ the double integral occurs in
\cite[Appendix~1, (7)--(8)]{TW}.

On the integration circles each $H$ is bounded and
$|1-w_iw_j|\ge 1-r^2$. Lemma~\ref{lem:analytic-det} with the
root values $w_i$ and angular values $v_i$ gives
\begin{equation}\label{eq:interior-site}
 |J_N(a,b;A,S,T)|\le\frac{C^N}{N!}e^{-\kappa N^2}r^{N(a+b)} .
\end{equation}
For real positive noncritical couplings $A>S+T$. The $v_i$ integral
has the sole interior pole $v_i=\rho_i$, with residue
$\rho_i^a/(S\sinh\gamma_i)$ in its one-body factor; $1/D$ cancels
the pole at zero. The roots have a uniform gap on the $w_i$ circles.
After taking the residues, \eqref{eq:split-pair} fills the apparent
reflected-pair poles as those circles expand to radius one.
Thus multiplying $J_N$ by $p^{N/2}$ gives the site integral in
\eqref{eq:real-site}, by reflection in the second coordinate.

Set $F_N=\sum_{a,b\in\Z}J_N(|a|,|b|;A,S,T)$. The parity of the
physical expansion gives integer powers in
\begin{align}
 \mathcal X_H&=(1-p^2)^{1/4}\sum_{N\ {\rm odd}}p^{(N-1)/2}F_N,
                  \label{eq:global-seriesH}\\
 \mathcal X_L&=(1-p^{-2})^{1/4}\sum_{N\ge2,\ N\ {\rm even}}p^{N/2}F_N.
                  \label{eq:global-series}
\end{align}
The integrals are holomorphic also when $ST=0$. Estimate
\eqref{eq:interior-site} on larger compact neighborhoods gives normal
convergence after every fixed mixed derivative.

\subsection{Identifying the state and summing over the lattice}\label{sec:lattice-sum}

For real positive noncritical couplings the sum in \eqref{eq:real-site}
thus gives $|C(a,b)-M^2|\le C'e^{-c(|a|+|b|)}$. With
$m_L=|\Lambda_L|^{-1}\sum_{x\in\Lambda_L}\sigma_x$ on centered
square boxes, averaging gives $\mu(\sigma_0m_L)\to M^2$.

In $L^2(\mu)$, box averages converge to projection onto the
translation-invariant functions. They fix that subspace and tend
to zero on differences $g-g\circ T$ for either coordinate shift
$T$, by telescoping. The closed span of these differences is its
orthogonal complement, and the averages are contractions. Write
$m$ for the projection of $\sigma_0$. Then
$\mu(m^2)=\mu(\sigma_0m)=M^2$, and
$\mu(\sigma_x\mid m)=m$ for every $x$ by invariance and projection.
Ignoring any fixed finite set of spins changes $m_L$ by $o(1)$.
An almost surely convergent subsequence therefore gives a
tail-measurable representative of $m$.

Conditioning on any positive-probability event determined by $m$
thus preserves the Gibbs conditional laws. Every such state lies
between the minus and plus states in ferromagnetic order. Indeed,
conditioning outside a finite box bounds increasing local observables
between the two fixed boundary conditions, by monotonicity of the
single-spin update in its neighbors. Exhaustion gives the infinite-volume order.
Its one-site mean lies in $[-M,M]$. Since the conditional mean is
$m$, we have $|m|\le M$ almost surely, hence $|m|=M$ by the second
moment. On $\{m=M\}$ the state is dominated by the plus state and
has the same mean at every site. For any real local observable $f$
supported on $V$, both $L\sum_{x\in V}\sigma_x\pm f$ are increasing
for large $L$. Applying the order to these observables gives equality
of the two expectations of $f$. On $\{m=-M\}$ the reversed comparison
gives the minus state. When $M=0$ the same comparison identifies the
unique state. Spin flip identifies the two-spin expectations in
either sign.

By \eqref{eq:interior-site} the two site sums and the particle sum
converge absolutely. At fixed $N$ sum first in $a$, obtaining
$(1+Z)/(1-Z)$. Summation in $b$ evaluates at zero the absolutely
convergent Fourier series of the density of the total angle. This
density at $\theta$ is the integral over $q_1,\ldots,q_{N-1}$ with
$q_N=\theta-\sum_{i<N}q_i$ and measure
$\prod_{i<N}dq_i/(2\pi)$. At zero it is exactly the integral in
\eqref{eq:complete-angular}, proving \eqref{eq:physical-formsH}--\eqref{eq:physical-forms}.
Equality on a real rectangle fixes the physical germ by the one-variable identity
theorem in each coordinate. Locally where $ST\ne0$, \eqref{eq:complete-angular} uses the
compatible $\sqrt B=\sqrt p/S$ with the corresponding HT factor
in \eqref{eq:physical-formsH}.

\subsection{Isotropic normalization and low-temperature descent}

At $z=w=t$ put $s=\sinh t$. Then $B=1$ and $C=s+s^{-1}$.
Define the single-variable coefficient on $|\zeta|>1$ by
\begin{equation}
 \begin{gathered}
 C_N(\zeta)=\frac1{N!}\int_{\sum q_i=0\ (2\pi)}
 \frac{1+Z}{1-Z}\prod_i\frac1{\sinh\gamma_i}
 \prod_{i<j}\frac{\sin^2((q_i-q_j)/2)}
                  {\sinh^2((\gamma_i+\gamma_j)/2)}
                  \prod_{i<N}\frac{\dd q_i}{2\pi},\\
 \cosh\gamma_i=\zeta+\zeta^{-1}-\cos q_i .
 \end{gathered}
 \label{eq:iso-coefficient}
\end{equation}
Choose each root $\rho_i=e^{-\gamma_i}$ with $|\rho_i|<1$.
On both positive real phases the roots and one-particle factors agree
with \cite[(3.17)--(3.18)]{ONGP}. The image
$\zeta+\zeta^{-1}$ avoids $[-2,2]$ on $|\zeta|>1$, so this
definition is single valued there and agrees with the complete
physical coefficient on either real phase. For these isotropic
functions we also write $\mathcal X_H(s)=X_H(s)$ and
$\mathcal X_L(s)=X_L(s)$. Equations
\eqref{eq:physical-formsH}--\eqref{eq:physical-forms} give
\begin{align}
 X_H(s)&=(1-s^4)^{1/4}s^{-1}
                   \sum_{N\ge1,\ N\ {\rm odd}}C_N(1/s),
                   \label{eq:iso-physicalH}\\
 X_L(s)&=(1-s^{-4})^{1/4}
                   \sum_{N\ge2,\ N\ {\rm even}}C_N(s).
 \label{eq:iso-physical}
\end{align}
The branches are the binomial ones. The summable bound on the angular
integrals gives normal convergence on compact subsets of each physical domain.

Under $t\mapsto t+2\pi i$, the quantities $S,T,A,p$, their canonical
roots and the integer-power coefficients in \eqref{eq:global-series}
are unchanged. The reduced LT sum is therefore invariant term by
term and descends under the covering of Lemma~\ref{geo:new:cover}.
The factor $\beta=t/(2J_x)$ remains on the temperature cover.

For $|\zeta|$ sufficiently large, the canonical root satisfies
$|\rho|\le C/|\zeta|$, the one-body factor is at most
$C/|\zeta|$, and every complete pair is at most $C/|\zeta|^2$.
The product kernel is uniformly bounded. Enlarging $C$ once gives
$|C_N(\zeta)|\le (C/|\zeta|)^{N^2}/N!$.
This supplies normal convergence also at $1/\zeta=0$.
At $\zeta=\infty$, direct expansion gives
$C_N(\zeta)=O(\zeta^{-N^2})$ for each fixed $N$ and
$C_1(\zeta)=1/(\zeta+\zeta^{-1}-2)$.
The joint construction is holomorphic at $s=0$; its Taylor
coefficients can be taken termwise on a sufficiently small disk.
Thus
\[
 X_H(s)=\frac{(1-s^4)^{1/4}}{(1-s)^2}+O(s^8),\qquad
 X_H(0)=1,\qquad
 X_L(s)=\frac1{4s^4}+O(s^{-6})\quad(s\to\infty).
\]
For even $N$, shifting every angle by $\pi$ and changing every
root's sign gives $C_N(-\zeta)=C_N(\zeta)$; conjugation follows
from the real germ. This justifies the even powers in the LT
expansion and will transfer its boundary obstructions between
quadrants. The LT leading coefficient follows already at $N=2$ from
$\tfrac12\int_0^{2\pi}\sin^2q\,\dd q/(2\pi)=1/4$.
For $N=1$ the root equation gives
$F_1=2\rho_1/[S(1-\rho_1)^2]=1/(A-T-S)$ by continuation. At $p=0$ the higher HT
terms vanish. On a decoupled axis this agrees with the independent
chains, which give
$\mathcal X_H(K_x,0)=\sum_{m\in\Z}\tanh(K_x)^{|m|}
=e^{2K_x}$ on the real axis and then by analytic continuation.
This decoupled restriction extends to an entire function.

\subsection{The energy-field response}
Finite-volume differentiation with respect to an energy field $h$
gives $\chi=\beta\mathcal X$. To pass this identity to the
specified infinite-volume state, we compare finite plus boxes using
field monotonicity. The GHS inequality for arbitrary
nonnegative interactions and site fields \cite[Theorem 1.5]{CGN}
says that a two-spin covariance decreases when any nonnegative
field increases. Pin the spins outside a smaller plus box with
fields tending to $+\infty$ and then release them to the fields
of a larger plus box. Covariances increase, so their nonnegative
sums increase to the finite connected sum just constructed.
Concavity in a uniform positive field gives
\[
 m_V(h)-m_V(0)\le h m'_V(0)\le\beta h\mathcal X.
\]
Taking the plus-state limit proves $m'(0+)\le\beta\mathcal X$.
For the reverse inequality turn on the field only on a finite set
$E$. This is a bounded local change of measure, with derivative
$\beta\sum_{v\in E}\operatorname{Cov}_{+,0}(\sigma_0,\sigma_v)$.
Ferromagnetic monotonicity bounds this magnetization by the
uniform-field one, with equality at $h=0$. Increase $E$ to the
lattice. Thus $m'(0+)=\beta\mathcal X$. Spin-flip symmetry makes
the HT derivative two-sided; the LT derivative is the plus-phase
right derivative. On a physical temperature slice, $\beta=t/(2J_x)$.


\section{Collision estimates at noncollinear edges}
\label{edge:new:section}

The product kernel can lose its boundary gap only when all angles
approach the endpoint.  The quadratic zeros of the pair factors control
this collision, while the remaining pairs provide decay in the particle
index. Both bounds hold on the constrained cycle with its $N-1$
independent angles.

Let $I\Subset(-1,1)$ be a compact interval, let
$\mathcal P\Subset\mathcal P^+$ be parameter polydisks, and suppose
$B(h,\eta)$ is bounded and holomorphic on a fixed neighborhood of
$I\times\overline{\mathcal P^+}$.  Fix $\sigma\in\{1,-1\}$ and assume
$\sigma\Im B(b,\eta)\ge\beta>0$ for $b\in I$ and
$\eta\in\mathcal P^+$.  Put $q_0=0$ or $\pi$ according as
$\sigma=1$ or $-1$, and write
\[
 W(h,q)=h+B(h)(\sigma-\cos q),\qquad
 \gamma(h,q)=\operatorname{arcosh}W(h,q),\qquad
 \gamma(b,q_0)=\ii\arccos b.
\]
On the physical side $h=b+\ii y$, $y>0$, write $\rho=e^{-\gamma}$ and
$Z=\prod_i\rho_i$.  The complete coefficient is
\[
 C_N(h)=\frac1{N!}\int_{\sum q_i=0\ (2\pi)}G_N(h,q)
              \frac{1+Z}{1-Z}\prod_{i<N}\frac{\dd q_i}{2\pi},
\quad
 G_N=B^{N^2/2}\prod_i\frac1{\sinh\gamma_i}
       \prod_{i<l}\frac{\sin^2((q_i-q_l)/2)}
       {\sinh^2((\gamma_i+\gamma_l)/2)}.
\]
All powers and branches are those fixed by the physical continuation
in Section~\ref{sec:physical}.

\begin{lemma}[A retained collision block]\label{edge:new:gaussian}
For each fixed integer $\ell\ge0$ let
$V_\ell=\prod_{a<b\le\ell}\sin^2((q_a-q_b)/2)$, with an empty
product equal to one.  There are $\kappa>0$ and common neighborhoods
such that $G_N=V_\ell A_{N,\ell}$, where the quotient is holomorphic
after its apparent coincidence singularities are filled, and
\[
 |\partial_h^j A_{N,\ell}|
 \le C_{\ell,j}\exp\{-\kappa N^2+C_\ell N\log(N+1)\},
 \qquad N\ge\ell.
\]
The bound is uniform in the real angles, in $b\in I$, and in the
larger parameter polydisk.  Every fixed additional parameter derivative
satisfies the same bound on $\mathcal P$.
\end{lemma}
\begin{proof}
On the compact branch $\Im\rho\le-d<0$ and $|\rho|\le1$.
Thus $1-\rho_i\rho_l$ and the one-body denominators are separated
from zero on a common complex neighborhood of the angular torus
and parameter set. The divided pair in \eqref{eq:split-pair} is uniformly bounded
there, including its anisotropy factor $B$. Apply
Lemma~\ref{lem:analytic-det} to the complete pair product, retaining
the $\ell$ labels. The remaining one-body factors, including
$B^{N/2}$, cost at most $C^N$.

This proves the estimate even
at coincident angles. Cauchy's formula on the larger parameter
set gives all the stated fixed derivatives.
\end{proof}

The numerator estimate supplies summability independently of the product
kernel.  Retaining a suitable finite block of collision zeros then makes
every prescribed derivative of that kernel integrable.

\begin{theorem}[Tail derivatives at a noncollinear edge]
\label{edge:new:budget}
Fix a parity, an integer cutoff $m\ge2$ of that parity, and $J\ge0$
such that
\[
                         2J<(m+2)^2-3.
\]
For $N\ge m+2$ of that parity and $0\le j\le J$,
\[
 \sup_{b\in I,\,0\le y\le y_0,\,\eta\in\mathcal P}
 |\partial_h^j\partial_\eta^\alpha C_N(b+\ii y,\eta)|
 \le\frac{C_{m,J,\alpha}N^J}{N!}
       e^{-\kappa N^2+C_mN\log(N+1)}.
\]
The boundary values are continuous for every $b\in I$, and the complete
tail and its derivative series converge uniformly on the closed
half-neighborhood. Each lower coefficient is
holomorphic through every subarc on which its endpoint resonance is
absent.  Consequently deleting the finite set
\[
 \{b\in I:N\arccos b\in2\pi\mathbb Z
             \text{ for some retained-parity }N\le m\}
\]
gives continuous jets through $J$ for the complete parity sum.
On a shorter collar let $j$ be the greatest integer with
$2j<n^2-3$, $n=m+2$. The order-$j$ derivative of the tail has modulus $O(\sqrt r)$ for even $n$
and $O(r\log(e/r))$ for odd $n$, where $r$ is distance in $h$.
Its lower derivatives are Lipschitz there. For $\sigma=-1$ in odd
parity the complete sum has continuous jets of every finite order
throughout $I$.
\end{theorem}
\begin{proof}
Write $q_i=q_0+x_i$ with $-\pi\le x_i\le\pi$. The compact branch
has bounded derivatives and
\[
 \Im W=y+\sigma\Im B(h)(1-\cos x_i)
       =\sinh(\Re\gamma_i)\sin(\Im\gamma_i).
\]
Hence $\Re\gamma_i\ge c(y+x_i^2)$. In the open collar expand
$(1+Z)/(1-Z)=1+2\sum_{L\ge1} Z^L$ and let $U_{L,N}$
denote the integral with factor $Z^L$, including $1/N!$.
For $L=0$ the integral extends to a fixed parameter disk with
the bound of Lemma~\ref{edge:new:gaussian}.
For $L\ge1$ retain the pair zeros on $n=m+2$ labels.
Lemma~\ref{edge:new:gaussian} bounds the divided numerator on a common
complex neighborhood, uniformly in the auxiliary parameters. With $n$
fixed, decreasing $\kappa$ absorbs the factor $\exp(C_n N\log(N+1))$
in its bound.
Proposition~\ref{p:transfer} in Appendix~\ref{app:forms}, with
$\vartheta=0$ and $d_N=\dist(Nq_0,2\pi\Z)$, gives
\eqref{eq:edge-transfer}. Cauchy's formula in the auxiliary polydisk
gives the same bound for each fixed $\partial_\eta^\alpha$.
Summation gives the stated moduli, and one more bounded derivative
gives Lipschitz continuity at lower orders. Uniform limits along
horizontal segments identify the tangential derivatives of the traces.

For each fixed lower index a pole at the endpoint requires all
$x_i=0$ and $N\arccos b\in2\pi\Z$. Else the compact contour has a
gap on a parameter disk. At the negative endpoint in odd parity
$d_N=\pi$, so $Q\ge\pi^2/N$ and
$\Re\sum_i\gamma_i\ge c(Ny+\pi^2/N)$.
A perturbation of $h$ by $c_1(y+N^{-2})$ keeps this positive,
since the sum changes by at most $CN|\Delta h|$. The kernel is then
$O(N)$. Cauchy's formula bounds its contribution after $j$ derivatives
by $j!(CN^2)^j CN e^{-\kappa'N^2}$, summably in odd $N$.
This proves the assertions on the entire negative arc.
\end{proof}

At a first finite resonance, take $m$ to be its index and $J$ to be
the order of its singular derivative. Theorem~\ref{edge:new:budget}
then bounds the derivatives of the entire higher tail. At a
collinear edge the same range of derivative orders is available.

\section{The collinear collar from lattice Fourier coefficients}
\label{iso:new:section}

We group directions on each square lattice shell. In an interior group
the Fej\'er kernel is concentrated near multiples of $2\pi$ in total angle.
Gaussian integrals on the contour for the central direction bound the
contribution of that group. At the endpoints we sum the individual site bounds.

For $\sigma=\pm1$ put $a_+=0$, $a_-=\pi$, and write
\[
 H_N^\sigma=\{\theta\in(\R/2\pi\mathbb Z)^N:
                    \textstyle\sum_i\theta_i=a_\sigma\},\qquad
 d\mu_N=\prod_{i<N}\frac{d\theta_i}{2\pi}.
\]
Eliminating the last coordinate gives $d\mu_N$ as hypersurface measure
divided by $\sqrt N(2\pi)^{N-1}$. For $\Im E>0$
choose $\rho_i=e^{-\gamma_i}$ by $|\rho_i|<1$ and
$\cosh\gamma_i=E-\cos\theta_i$.  Define
\begin{equation}\label{iso:new:F}
 \mathcal F_N^\sigma(E)=\frac1{N!}\int_{H_N^\sigma}
 \frac{1+\sigma Z}{1-\sigma Z}
 \prod_i\frac1{\sinh\gamma_i}\prod_{i<l}P_{il}\,d\mu_N,
 \quad Z=\prod_i\rho_i,
 \quad P_{il}=\frac{\sin^2((\theta_i-\theta_l)/2)}
                       {\sinh^2((\gamma_i+\gamma_l)/2)}.
\end{equation}
The same sign appears in the momentum constraint and the product
kernel. The physical coefficient uses $\sigma=+$, and $\sigma=-$
implements the dispersion reflection for odd orders.

\begin{theorem}[Tail derivatives at an isotropic edge]\label{iso:new:budget}
Let $K\Subset(0,2)$ be a compact interval, fix an integer $n\ge2$,
and let $j_*$ be the greatest integer satisfying $2j_*<n^2-3$.
There exist $y_0,\kappa>0$ such that for either sign, every $N\ge n$
and $0\le j\le j_*$, the derivative $\partial_E^j\mathcal F_N^\sigma$
extends continuously to $K+\ii[0,y_0]$ and satisfies
\begin{equation}\label{iso:new:majorant}
 \sup_{E\in K+\ii[0,y_0]}|\partial_E^j\mathcal F_N^\sigma(E)|
       \le \frac{C_j^N}{N!}e^{-\kappa N^2}.
\end{equation}
The constants may depend on $K,n,j$. For the sum over $N\ge n$,
or its subseries of either parity, the derivative of order $j_*$
has modulus of continuity
\[
  O(h^{1/2})\quad(n\ \text{even}),\qquad
  O\bigl(h\log(e/h)\bigr)\quad(n\ \text{odd})
\]
on shorter closed collars, with $h$ the distance between two points.
Every lower derivative is Lipschitz there. The same statements hold
for compact intervals in $(-2,0)$ by the signed reflection
\eqref{eq:iso-reflection}.
\end{theorem}

\subsection{The two lattice coordinates}\label{sec:site-saddle}

Put $\gamma=\ii\phi$ and $z=e^{\ii q}$. For $A\in\Z_{\ge0}$ and $B\in\Z$ let
\[
 I_N(A,B;E)=\frac1{N!}\int_{\T^N}
  \prod_{i<j}P_{ij}\prod_i
  e^{\ii(Bq_i-A\phi_i)}\frac{\dd q_i}{2\pi\sinh\gamma_i}.
\]
Reflection of all angles gives $I_N(A,-B)=I_N(A,B)$.
The residue calculation in \eqref{eq:double-site} with spectral
parameters $(E,1,1)$ gives $I_N(A,B)=J_N(A,|B|;E,1,1)$,
first for real $E>2$ and then on the slit plane. The two circles
therefore give $I_N(A,B)=I_N(B,A)$ for $A,B\in\Z_{\ge0}$ in both parities.

Expanding the geometric kernel of \eqref{iso:new:F} in the first
lattice coordinate and evaluating the Fourier series of the total
angle at $a_\sigma$ in the second gives
$\mathcal F_N^\sigma=\sum_{a,B\in\Z}\sigma^{a+B}I_N(|a|,B)$.
Both lattice sums converge absolutely by \eqref{eq:interior-site}
for $(E,1,1)$.
Group sites by $A=\max(|a|,|B|)$. Reflections and axis exchange give
\begin{equation}\label{eq:signed-shell}
 \mathcal F_N^\sigma=I_N(0,0)+8\sum_{A\ge1}\sigma^A S_{A,N}^\sigma,
 \qquad S_{A,N}^\sigma=\sum_{B=0}^A{}'\,\sigma^B I_N(A,B).
\end{equation}
The prime gives half weight to $B=0,A$, whose symmetry orbits have
four sites each. Every interior orbit has eight.

\subsection{Angular contours and the sum over lattice directions}\label{sec:direction-sum}

Fix $x\in(0,2)$, put $b=\arccos(x-1)$, and set
$g(r)=\arccos(x-\cos r)$ on $(-b,b)$. On the upper boundary
the physical lift there is $\phi=g$. We have
\[
 g'=-\sin r/\sin g,\qquad
 g''=-x(1-\cos r\cos g)/\sin^3g<0 .
\]
Let $f$ be odd and $2\pi$-periodic, equal to $r$ on $[0,b]$ and
linear from $b$ to $0$ on $[b,\pi]$. It has unique extrema at $\pm b$.

\begin{lemma}\label{lem:site-contour}
Let $c$ range in a compact subinterval of $(-b,b)$. At $E=x+\ii y$
write $q=r+\ii v(f(r)-f(c))$, $-\pi\le r\le\pi$.
For $0\le v\le v_0$ and $0\le y\le y_0$ the lifted angle satisfies
\begin{equation}\label{eq:site-phase}
 \Im(g'(c)q-\phi)\ge a\bigl(y+v(r-c)^2\bigr).
\end{equation}
For $v+y>0$ the one-body and pair denominators do not vanish in the
homotopy. At $v=v_0$ they have uniform gaps on a complex neighborhood
of the graph and the parameter values. The constants are locally
uniform in $x$ for centers a positive distance from the endpoints.
\end{lemma}
\begin{proof}
On a closed interval inside $(-b,b)$ containing the centers, put
$h=f(r)-f(c)=r-c$. At $y=0$ the imaginary part is
$v h(g'(c)-g'(r))+O(v^3|h|^3)$, whose leading term is at least
$a_1v(r-c)^2$ by strict concavity. Increasing $y$ adds at least
$a_2y$ since $\Re(1/\sin\phi)>0$ there.

Near either endpoint $\sin r\,h\ge\delta>0$, hence
$\Im W\ge y+\delta v$ for $W=E-\cos q$. Here the lift has
$G=-\Im\phi>0$ when $\Im W>0$, and
$\Im W/G=\sin(\Re\phi)\sinh(G)/G\to0$ as $W\to1$.
Taking the endpoint neighborhoods small, then taking $v_0,y_0$ small,
makes $G/(y+\delta v)$ large enough to absorb $v|g'(c)h|$ and prove
\eqref{eq:site-phase}. On the remaining compact part of $|r|>b$,
$G$ is already positive at $v=y=0$. Use the interior estimate up
to the endpoint neighborhoods.

Every unit root on the original cycle at $y=0$ has argument in
$[-b,0]$, where $b<\pi$. A pair product can equal one only if both
roots are the branch value $1$. In the endpoint neighborhoods $\Im W>0$
puts the roots strictly inside the unit disk throughout the homotopy
for $v+y>0$. Compactness gives the gaps elsewhere and the stated
complex neighborhoods of the final graphs.
\end{proof}

For a direction $0\le B\le A$, choose $g'(c)=B/A$. Then
$-\arccos(x/2)\le c\le0$. On the final contour of
Lemma~\ref{lem:site-contour} the $N$ transfer and Fourier factors have
modulus at most $\exp[-aA(Ny+\sum_i(r_i-c)^2)]$.
Retain the pairs on $n$ labels in Lemma~\ref{lem:analytic-det}.
The divided numerator is bounded by $C_n^N e^{-\kappa N^2}/N!$
including the one-body factors, the contour Jacobian and the factorial.
Since each angle follows the same Lipschitz graph, the retained
differences cost at most $C_n\Delta_n(r)^2$, where
$\Delta_n(r)=\prod_{1\le i<k\le n}(r_i-r_k)$.
With $d=[N+n(n-1)]/2$, integration of the $N$ Gaussians gives
$C_n^N A^{-d}$: put $r_i=c+A^{-1/2}u_i$, and integrate the fixed
polynomial $\Delta_n(u)^2$ against the product Gaussian.
On an $E$-disk of radius $\epsilon/(1+A)$ the phases change by
$O((1+A)^{-1})$ per angle, at a cost $C^N$ in the exponential.
These bounds also hold for a group of directions on its central contour,
with the Fourier factor below.

Put $L=\lceil\sqrt A\rceil$ and $w_L(b)=(1-|b|/L)_+$.
The translates $w_L(B-kL)$ sum to one. The at most four groups meeting
an endpoint of $[0,A]$ contribute at most $O(L)$ times the single-site
bound. In every other nonempty group put $m=kL$, choose $g'(c_m)=m/A$
and deform all its terms to that one contour. Apart from the central
exponential and the sign $\sigma^m$, the sum inserts
\[
 W_L(T)=\sum_b w_L(b)e^{\ii bT}
 =\frac1L\left(\frac{\sin(LT/2)}{\sin(T/2)}\right)^2,\qquad
 T=a_\sigma+\sum_i q_i .
\]
The finite sum and the sine formula give
$|W_L(T)|\le CL e^{L|\Im T|}/(1+L\dist(\Re T,2\pi\Z))^2$.
Write $U=\sum_i|r_i-c_m|$. Then $|\Im T|\le C U$ and
\[
 1+L\dist(Nc_m+a_\sigma,2\pi\Z)
 \le(1+LU)(1+L\dist(\Re T,2\pi\Z)).
\]
Half the Gaussian in $\sum_i(r_i-c_m)^2$ absorbs
$e^{CLU}(1+LU)^2$ at a cost $C^N$, since $L^2/A\le4$.
The group therefore costs the integrated single-site bound times
$CL/(1+L\dist(Nc_m+a_\sigma,2\pi\Z))^2$.

Consecutive centers $c_m$ are separated by at least $cL/A$,
since $|g''|$ is bounded on their interval. Group the numbers
$Nc_m+a_\sigma$ by their nearest multiple of $2\pi$. Only $O(N+1)$
multiples occur. Within each group, after subtracting that
multiple and scaling by $L$, the numbers are separated by at least
$cNL^2/A\ge c$. The sum of $(1+|v|)^{-2}$ over such a set is bounded.
Thus summing the directions costs $O((N+1)L)$.
Cauchy's estimate on the same $E$-disk gives
\begin{equation}\label{eq:iso-shell-est}
 |\partial_E^j S_{A,N}^\sigma(x+\ii y)|
 \le\frac{C_{n,j}^N}{N!}e^{-\kappa N^2}
     A^{j-[N-1+n(n-1)]/2}e^{-a'Ay}.
\end{equation}
The site at the origin has the same particle bound on a fixed disk
by the graph with $c=0$.

\begin{proof}[Conclusion of Theorem~\ref{iso:new:budget}]
The power in \eqref{eq:iso-shell-est} is
$j-\alpha_{N,n}$ as in \eqref{eq:edge-transfer}. Summation thus
gives \eqref{iso:new:majorant} and both continuity moduli by the same
split at $A=h^{-1}$. Equation~\eqref{eq:iso-reflection} below gives
the negative intervals.
\end{proof}

\subsection{The complement of the finite collision charts}\label{sec:finite-shift}

At fixed $N\ge2$, constant branch configurations and regular constant
resonances can occur only at the respective values
\begin{equation}\label{iso:new:exceptions}
 1+\cos\frac{2\pi l+a_\sigma}{N},\qquad
 \cos\frac{2\pi l+a_\sigma}{N}+
 \cos\frac{2\pi k+a_\sigma}{N},\qquad 0\le l,k<N .
\end{equation}
Suppose $x$ is not one of the branch values. On $H_N^\sigma$ take
$c=N^{-1}\sum_i f(r_i)$. No tuple has all its angles at the same
endpoint $\pm b$, so the unique extrema of $f$ and compactness
keep these centers in a compact subinterval of $(-b,b)$. The
displacements sum to zero and preserve the signed constraint.
Equation \eqref{eq:site-phase} gives
$-\sum\Im\phi_i\ge a(Ny+v\sum_i(r_i-c)^2)$. At $v>0$ a vanishing
product denominator at the boundary therefore requires
$r_i=c=q_0$ for all $i$ and both $Nq_0,Ng(q_0)\equiv a_\sigma$
modulo $2\pi$.

Leave a small real chart at each such tuple by multiplying $v$ by
a common smooth cutoff, zero near the center and one outside the
charts. On a punctured real chart $\sum_i g(r_i)<Ng(q_0)$ by
strict concavity and $\sum_i(r_i-q_0)=0$. Choose the charts small
enough to exclude all other product phases. The transition annuli
thus have a gap also where the displacement vanishes. Deform at
$y>0$ and insert cutoffs supported in the retained real charts,
equal to one near their centers. The complement has uniform gaps
in all denominators on the final cycle and extends holomorphically
to an $E$-disk about $x$. If there are no regular constant resonances,
this continues the entire coefficient. For $N=1$ the two signed
coefficients are $1/(E-2)$ and $1/(E+2)$.

Define the lower-half-plane branch by complex conjugation of the
real germ. On these two branches the coefficient identities
\[
 \mathcal F_N^\sigma(-E)=(-1)^N
                         \mathcal F_N^{(-1)^N\sigma}(E)
\]
follow by replacing $\theta_i$ by $\theta_i+\pi$ and $\rho_i$ by
$-\rho_i$. Each pair is unchanged, every one-body factor
changes sign, and both products acquire $(-1)^N$.
Equivalently, using only values with $\Im E>0$,
\begin{equation}\label{eq:iso-reflection}
 \mathcal F_N^\sigma(E)=(-1)^N
 \overline{\mathcal F_N^{(-1)^N\sigma}(-\overline E)}.
\end{equation}
For odd $N$ the signed identities transport the estimates to negative
dispersion by changing the sign in both the constraint and the kernel.

\begin{corollary}[Complete finite jets]\label{iso:new:jets}
For each $r\ge0$ and either parity, choose $m\ge2$ in that parity
with $2(r+1)<(m+2)^2-3$. On a compact interval in $(0,2)$ avoiding
\eqref{iso:new:exceptions} for the chosen parity through $N=m$, the complete parity sum has continuous
derivatives through order $r+1$ from its physical side. Its derivatives
through order $r$ are Lipschitz on shorter closed collars. The reflected
statement in $(-2,0)$ uses the exceptions for the sign in
\eqref{eq:iso-reflection}. The corresponding physical factors
$(1-\zeta^{-4})^{1/4}$ and
$\zeta(1-\zeta^{-4})^{1/4}$, and the noncritical map
$E=\zeta+\zeta^{-1}$, preserve this conclusion.
\end{corollary}
\begin{proof}
With first tail index $m+2$, Theorem~\ref{iso:new:budget} gives
uniform convergence through order $r+1$ on the closed collar.

The finitely many lower coefficients extend holomorphically across
the interval. Integration along a segment gives the Lipschitz bounds.
Passing to the uniform boundary limit in
\[
 F^{(j)}(x'+iy)-F^{(j)}(x+iy)
   =\int_x^{x'}F^{(j+1)}(u+iy)\,du
\]
identifies tangential derivatives of the traces. The chain and product
rules give the same regularity in the physical coordinate.
\end{proof}

\section{Constrained periods and noncollinear boundary singularities}
\label{sec:period}

At a positive endpoint the first resonant coefficient has a nonzero
leading period. Its evaluation uses a constrained Gaussian integral
and a beta integral, with the independent-coordinate measure of the
physical coefficient. The same periods enter the radial asymptotics in
Section~\ref{iso:new:radial:section} and the residue calculation in
Section~\ref{neg:new:section}.

Lemma~\ref{lem:local-period} in Appendix~\ref{app:forms} evaluates the
local integral with this measure. We apply it first at a noncollinear edge.

\subsection{A nonzero finite collision period}

In the positive endpoint chart of Section~\ref{edge:new:section},
fix an odd prime $p$, $1\le r\le(p-1)/2$, and set
\[
 \phi=2\pi r/p,\quad b=\cos\phi,\quad
 h=b+u,\quad B_0=B(b),\quad \Im B_0>0.
\]
Along a transverse curve, $B=B(b+u)$ is holomorphic in $u$. Let $n=p$
for the odd sum and $n=2p$ for the nonzero-even sum; write
$a_n=(n^2-1)/2$, $k=\lfloor a_n\rfloor$ and
$\nu=k+1-a_n$, so $\nu=1$ or $1/2$ respectively. The Gaussian
factor $P_n=\prod_{j=1}^{n-1}(j!/2^j)$ is as in Appendix~\ref{app:forms}.

\begin{proposition}[Full-tail noncancellation at a positive edge]
\label{prop:positive-edge}
For the complete coefficient $C_n$ and $u=i\epsilon$,
$\epsilon\downarrow0$,
\begin{equation}
 \lim u^\nu\partial_u^k C_n(b+u)=L_{n,B},\qquad
 L_{n,B}=
 \frac{2^{1-a_n}i^{1-n^2}(-1)^k\sqrt{B_0}\,P_n
         n^{a_n-3/2}\Gamma(\nu)}
      {\pi^{(n-1)/2}\sin^{2a_n}\phi}\ne0.
 \label{eq:positive-L}
\end{equation}
Here $u^{1/2}=e^{i\pi/4}\sqrt\epsilon$. For odd $n$, $\sqrt{B_0}$
is the local branch in the physical coefficient. For even
$n$, it is the beta-period branch with $\arg B_0\in(0,\pi)$,
equivalently $e^{i\pi/4}\sqrt{-iB_0}$ with the right-half-plane
square root. The even-particle coefficient has an integer power of
$B$. The beta period fixes the branch of the half-power in the leading
term.
All lower retained coefficients are
holomorphic near the boundary point; all higher coefficients together
have a bounded $k$th derivative. Consequently the full physical sum
is nonextendable on every such nonramified transverse curve.
The estimates and leading limit are locally uniform on the stipulated
compact holomorphic parameter families for a fixed selected phase.
\end{proposition}
\begin{proof}
On the real boundary, $|Z|=1$ is possible only when all angles
equal the endpoint zero. Its product resonance is $N\phi\in2\pi\Z$.
Thus $n$ is the first index of the retained parity and no lower
coefficient has a denominator zero. The canonical roots give a
holomorphic neighborhood for each of those finitely many coefficients.

Near the unique equality point use $q_n=-\sum_{i<n}q_i$,
$Q=\sum q_i^2$ and $\Delta=\Delta_n(q)$. The expansion is
\[
 \gamma(q)=i\phi-\frac{i}{\sin\phi}
              (u+B_0q^2/2)+O(|u|^2+|u|q^2+q^4).
\]
Take $D=i\sin\phi\,(\sum\gamma_i-in\phi)$. Each pair has an exact
double diagonal zero. After removal of $\Delta^2$, multiplication
of the integrand by $D$ fills the simple pole of the product kernel
and gives a symmetric analytic amplitude with central value
\[
 A_0=\frac{2i^{1-n^2} B_0^{n^2/2}}
 {n!(2\pi)^{n-1}2^{n(n-1)}\sin^{n^2-1}\phi}.
\]
In Lemma~\ref{lem:local-period} take $(d,h,\tau)=(n,B_0/2,i)$.
Both coefficients $ni,B_0/2$ lie in the open upper half-plane.
Its formula gives \eqref{eq:positive-L} after multiplication by
$u^\nu$. For odd $n$, the square root of $B_0$ comes from
$B_0^{n^2/2}$. For even $n$, the beta period continued in the upper
half-plane gives the stated root. Away from this chart the denominator
has a fixed gap for the first coefficient.

Apply Theorem~\ref{edge:new:budget} with $m=n$ and $J=k$.
It bounds the sum of all terms of higher particle number in the same parity.
The same local integral at lower derivative orders is bounded
in even parity and at most logarithmic in odd parity. Multiplication
by the physical prefactor therefore multiplies the leading coefficient
by the nonzero boundary value of that prefactor. A nonramified curve
uses $u$ as a local coordinate and has the same conclusion.

For compact holomorphic parameter families, reserve a larger polydisk
on which the
strict imaginary margin, branch gaps and nonzero physical factors
persist. All estimates above then have common constants. Cauchy's
formula in the auxiliary parameters gives every fixed mixed estimate.
For each family member the nonzero value in \eqref{eq:positive-L}
is the leading coefficient for the transverse derivative $\partial_u^k$.
\end{proof}

\subsection{The other orientations and the remaining odd edge}

In a regular edge chart, the physical side satisfies
$\operatorname{sign}(\Im h)=\operatorname{sign}(\sigma\Im B)$.
If this is negative, conjugate the chart and curve, using the real
physical germ; the conjugated holomorphic curve is
$\widetilde B(h)=\overline{B(\overline h)}$. This reduces to the
upper physical side.
At a negative endpoint with even $N$, shift every angle by $\pi$.
The constraint is unchanged and $B$ is replaced by $-B$, whose
imaginary part is positive. Also $N^2/2$ is an even integer, so
$B^{N^2/2}=(-B)^{N^2/2}$ exactly. Proposition~\ref{prop:positive-edge}
therefore proves the negative even obstruction with its full phase.
For odd $N$, the same shift changes the constraint to an odd winding,
which excludes endpoint equality at every finite index.
Section~\ref{neg:new:section} proves the full-sum obstruction from the
growth of the boundary jets.


\section{First resonances and sharp radial asymptotics}
\label{iso:new:radial:section}

On the isotropic prime families below, the first resonant coefficient
has exactly four singular configurations. Their complex amplitudes
are equal and add in the physical sum. Absolute convergence of the
higher particle tail at the resonant derivative order preserves the
leading amplitude and the radial remainder.

Let $p\ge11$ be prime, choose distinct
$a,b\in\{1,\ldots,(p-1)/2\}$, and set
\[
 \alpha=2\pi a/p,\quad\beta=2\pi b/p,\quad
 E_*=\cos\alpha+\cos\beta=2\cos\theta,\quad
                            \zeta_*=e^{i\theta},\quad0<\theta<\pi.
\]
For the even family require $\alpha,\beta\in(0,\pi/2)$ and
put $n=2p$.  For the odd family require both angles in
$(0,\pi/2)$ or both in $(\pi/2,\pi)$, and put $n=p$.
Write $C_N(\zeta)=\mathcal F_N^+(\zeta+\zeta^{-1})$.

\begin{lemma}[The first equality resonance]\label{iso:new:radial:arithmetic}
The first order of the specified parity satisfying
$E_*=\cos(2\pi l/N)+\cos(2\pi q/N)$ is $n$.
At that order the unordered cosine pair is uniquely
$\{\cos\alpha,\cos\beta\}$, so the constrained angular
coefficient has exactly four resonant equality points.
No endpoint equality is resonant at any order.  The even
family is dense in the upper-right exterior unit arc; the
two odd families are dense in both upper exterior quadrants.
\end{lemma}
\begin{proof}
Put $\omega=e^{2\pi i/p}$, $d=(p-1)/2$,
$u_j=\omega^j+\omega^{-j}$.
In the real cyclotomic field of degree $d$ the conjugates of each
$u_j$ sum to $-1$. Thus the algebraic integer $v=u_a+u_b$ has
trace $-2$, so it is nonzero and has a negative conjugate. If all
conjugates were negative, their product would have modulus at most
$(2/d)^d<1$ by the arithmetic--geometric mean inequality, whereas
its norm is a nonzero integer. Hence $v$ has both signs among its
conjugates and is irrational.

If $p\nmid N$, choose an exponent giving a different conjugate
of $v$. The Chinese remainder theorem lifts it modulo $pN$ to an
exponent congruent to one modulo $N$. The resulting automorphism
fixes every $N$th root and moves $v=2E_*$. This excludes every
smaller index of the prescribed parity.

The sign change excludes both $1+\cos r$ and $-1+\cos r$ for
any root-of-unity angle $r$. At index $p$ or $2p$ the remaining
cosines belong to $\{\pm u_j/2\}$. The $u_j$, $1\le j\le d$,
are rationally independent: a relation lifts to a polynomial of
degree at most $p-1$ with zero constant term divisible by
$1+X+\cdots+X^{p-1}$. Matching coefficients gives the unique
unordered pair, with constant configurations at $\pm\alpha,\pm\beta$.
Finally, two distinct adjacent prime-grid angles can approach
any chosen angle in either half interval.  Their averaged
cosines prove the stated density.
\end{proof}

Set
\[
 a_n=\frac{n^2-1}{2},\quad k=\lfloor a_n\rfloor,\quad
 \nu=k+1-a_n,\quad H=E_*(1-\cos\alpha\cos\beta),\quad
 P_n=\prod_{j=1}^{n-1}\frac{j!}{2^j}.
\]
Thus $\nu=1/2$ in the even family and $\nu=1$ in the odd
family. With $d\mu_n=(2\pi)^{1-n}\dd q_1\cdots\dd q_{n-1}$, the four
resonant configurations have the same nonzero amplitude
\begin{equation}\label{iso:new:radial:L}
 L_*=
 \frac{i^{-n^2}(-1)^k\Gamma(\nu)P_n n^{a_n-3/2}
          (\sin\theta)^{a_n-1}e^{-ik\theta}}
      {\pi^{(n-1)/2}(-iH)^{a_n}}.
\end{equation}
For even $n$, the last power has $\arg(-iH)=-\pi/2$.
For odd $n=p$, $a_n=k$ is divisible by four, so the
same formula applies to either sign of $H$.

\begin{proposition}[A sharp first-particle period]\label{iso:new:radial:first}
Along $\zeta_\epsilon=(1+\epsilon)\zeta_*$,
\[
 C_n^{(k)}(\zeta_\epsilon)=4L_*\epsilon^{-\nu}
       +\begin{cases}O(1),&n\text{ even},\\
                  O(1+|\log\epsilon|),&n\text{ odd}.
         \end{cases}
\]
Every fixed derivative of each smaller coefficient of that
parity, and of the complement of the four first-order
charts, is bounded.  Below order $k$, the first coefficient
has bounded derivatives in the even case, while the odd case
has bounded derivatives through $k-2$ and a logarithmic
$(k-1)$st derivative.
\end{proposition}
\begin{proof}
Use $z=\zeta/\zeta_*-1$ and one lower chart
$q_i=-\alpha+u_i$, $\sum_i u_i=0$, $Q_2=\sum_i u_i^2$.
At its center $\gamma_i=i\beta$. From $\cos q+\cos\phi=E$ one gets
$\phi_q=\sin\alpha/\sin\beta$ and
$\phi_{qq}=-H/\sin^3\beta$. The constraint cancels the linear term.
The denominator $D=1-Z$ has coefficients
\[
 d=\frac{2n\sin\theta}{\sin\beta},\qquad
 h=-\frac{iH}{2\sin^3\beta},\qquad D=d z+hQ_2+
            O(|z|^2+|z|Q_2+|u|^3).
\]
Here $d>0$ and $h$ is nonzero imaginary. The symmetric analytic
amplitude after removing $\Delta_n^2$ has central value
\[
 A_0=\frac{2i^{-n^2}}
 {n!(2\pi)^{n-1}2^{n(n-1)}\sin^{n^2}\beta}.
\]
Indeed the product numerator is two, each one-body term is
$(i\sin\beta)^{-1}$ and each divided pair contributes
$1/[4(i\sin\beta)^2]$. Apply Lemma~\ref{lem:local-period} with
$\tau=1$ and multiply by $\zeta_*^{-k}$ to return to the
$\zeta$ derivative. Substituting $G_n$ gives \eqref{iso:new:radial:L},
since $-n^2-(a_n-1)+3a_n=0$. This coefficient is symmetric in
$\alpha,\beta$. Exchanging them and reflecting all angles thus
gives the same complex number at each of the four points.

At positive dispersion Lemma~\ref{iso:new:radial:arithmetic}
excludes every full branch equality and identifies the four ordinary
resonant equalities. Section~\ref{sec:finite-shift} continues the
complement holomorphically, with the contour unchanged on those four
charts. The same argument continues every lower fixed coefficient.
At negative dispersion and odd particle number, the reflection
\eqref{eq:iso-reflection} changes the sign in both the momentum
constraint and the product kernel.
The lower-derivative bounds in the selected charts follow
from the same radial integral with a smaller pole power.
\end{proof}

To pass from the coefficient to the full susceptibility, restore the
physical factors on the upper exterior ray:
\[
 p_L(\zeta)=(1-\zeta^{-4})^{1/4},\quad
 p_H(\zeta)=\zeta(1-\zeta^{-4})^{1/4},\qquad
 \mathcal A_\eta=4\zeta_*^k p_\eta(\zeta_*)L_*\ne0.
\]
Let $g_L(\epsilon)=\mathcal X_L((1+\epsilon)\zeta_*)$
and $g_H(\epsilon)=\mathcal X_H(((1+\epsilon)\zeta_*)^{-1})$.

\begin{theorem}[Sharp complete physical radial remainders]
\label{iso:new:radial:full}
On the respective selected prime-angle families,
\begin{align*}
 g_L^{(k)}(\epsilon)&=\mathcal A_L\epsilon^{-1/2}+O(1),\\
 g_L(\epsilon)&=\sum_{j=0}^{k-1}\frac{b_j}{j!}\epsilon^j
 +\frac{\sqrt\pi\mathcal A_L}{\Gamma(k+1/2)}\epsilon^{k-1/2}
 +O(\epsilon^k),\\
 g_H^{(k)}(\epsilon)&=\mathcal A_H\epsilon^{-1}
                                   +O(1+|\log\epsilon|),\\
 g_H(\epsilon)&=\sum_{j=0}^{k-1}\frac{b_j}{j!}\epsilon^j
 +\frac{\mathcal A_H}{(k-1)!}\epsilon^{k-1}\log\epsilon
 +O(\epsilon^k\log(e/\epsilon)).
\end{align*}
In LT all $b_j$ are finite radial jets and
$g_L^{(k-1)}(\epsilon)-b_{k-1}
=2\mathcal A_L\epsilon^{1/2}+O(\epsilon)$.
This derivative is exactly H\"older $1/2$ at zero.
In HT, $b_0,\ldots,b_{k-2}$ are finite radial jets and
$b_{k-1}$ is the finite part after subtraction of the
displayed logarithmic monomial. Explicitly, if
$H_{k-1}=\sum_{l=1}^{k-1}l^{-1}$, then
$b_{k-1}=\lim_{\epsilon\downarrow0}
[g_H^{(k-1)}(\epsilon)-\mathcal A_H(\log\epsilon+H_{k-1})]$.
One has
$g_H^{(k-2)}(\epsilon)-b_{k-2}
=\mathcal A_H\epsilon\log\epsilon+O(\epsilon)$.
This derivative is log-Lipschitz and H\"older of every
exponent less than one, but is not Lipschitz; the next
derivative diverges logarithmically.
\end{theorem}
\begin{proof}
The higher tail starts at $n+2$. Since $2k\le n^2-1<(n+2)^2-3$,
Theorem~\ref{iso:new:budget} gives a continuous bounded $k$th
derivative up to the edge.

Identity~\eqref{eq:iso-reflection} gives the same bound on the
negative odd family.
Combining this full tail with
Proposition~\ref{iso:new:radial:first} proves the top
derivative asymptotics.  The analytic physical factors
are inserted before differentiation: their lower
product-rule terms are bounded in LT and at most
logarithmic in HT.  Since $d/d\epsilon=\zeta_*d/d\zeta$,
the coefficient is precisely $\mathcal A_\eta$.

For LT, subtract the displayed fractional power.  Its
$k$th derivative is exactly
$\mathcal A_L\epsilon^{-1/2}$, so the difference has
bounded $k$th derivative and Taylor remainder
$O(\epsilon^k)$.  Integrating once gives the asserted
top finite-jet increment; integrating its bound between
two radii proves H\"older $1/2$.  Its nonzero coefficient
excludes every larger exponent.
For HT, subtract the displayed logarithmic monomial.
Its $k$th derivative is exactly $\mathcal A_H/\epsilon$.
The remaining $k$th derivative is integrable and
logarithmically bounded, so integral Taylor's formula
gives the stated remainder and finite part.  Differentiating
the explicit monomial $k-2$ times gives the increment.
Integration of the logarithmic next derivative proves
the log-Lipschitz bound, and the nonzero coefficient
excludes Lipschitz continuity.
\end{proof}

These first resonances also attain the two moduli in
Theorem~\ref{iso:new:budget} for a tail starting at the resonant
index $n$. Its highest admissible order is $k-1$ for $n=2p$ and
$k-2$ for $n=p$. The lower particle coefficients are analytic there.
The physical prefactor is nonzero, and the radial derivative of
$E=\zeta+\zeta^{-1}$ at $\epsilon=0$ is $2\ii\sin\theta\ne0$.
Lower derivatives of the tail are Lipschitz, so their contributions
to the increments under these changes are $O(\epsilon)$. The nonzero
$\sqrt\epsilon$ and $\epsilon\log\epsilon$ terms therefore give the
same sharp orders for the physical parity tails.

Evenness and conjugation carry the even
family to the other quadrants.  The two odd families, together with
inversion and conjugation, cover a dense subset of the physical HT
circle. The resulting dense set of nonzero singular periods supplies
the local obstructions used in the componentwise boundary argument of
Section~\ref{sec:assembly}.

\subsection{The first diagonal resonance}

\begin{corollary}\label{iso:prime-power-arcs}
Let $\zeta$ be a primitive $d$th root. In HT take any odd $d\ge5$
and set $n=d$. In LT take any $d\ge5$ other than six and set
$n=\operatorname{lcm}(d,2)$. The first singular coefficient of the
physical parity has order $n$, with exactly two constant resonances.
With $k=\lfloor(n^2-1)/2\rfloor$, its $k$th exterior radial derivative,
and that of the full normalized susceptibility, have
nonzero leading terms of order $\epsilon^{-1}$ in HT and
$\epsilon^{-1/2}$ in LT. For HT the physical variable is $s=1/\zeta$
and the exterior radius is in $\zeta$; for LT it is in $s$ itself.

In particular, take $n=3^m$, $m\ge2$, in HT or $n=2^m$, $m\ge3$,
in LT. Every open arc of length $\ell>4\pi/n$ contains a point
where a real $C^k$ extension from the physical side is impossible.
In LT one may take $k\le\max(32,32\pi^2/\ell^2)-1$, and in HT
$k\le(\max(81,144\pi^2/\ell^2)-1)/2$. For either phase one can
in fact take $k\le(8\pi^2+o(1))\ell^{-2}$ as $\ell\downarrow0$.
\end{corollary}

\begin{proof}
Put $E_*=\zeta+\zeta^{-1}$. An identity $E_*=\pm1+\cos r$ at a root-of-unity angle
would make $\cos r$ an algebraic integer. Its conjugates lie in
$[-1,1]$, being cosines of root-of-unity angles. An endpoint conjugate
forces $\cos r=\pm1$. Otherwise its integer norm has modulus below one,
so $\cos r=0$. Then $E_*\in\{0,\pm1,\pm2\}$, whose root orders
are $1,2,3,4,6$, all excluded here.

Suppose two $N$-grid cosines sum to $E_*$. We may send $\zeta$ to
any primitive $d$th root while permuting the $N$-grid: lift its unit
exponent modulo $d$ to a unit modulo $\operatorname{lcm}(d,N)$ by
prescribing one at each additional prime divisor of $N$. Neither summand
is an endpoint, also after conjugation. The grids with $N\le2$ have
only endpoints.

When $n=d$, send $\zeta$ to $e^{2\pi\ii/d}$. The largest remaining
$N$-grid cosine is $\cos(2\pi/N)$. The target $2\cos(2\pi/d)$ is too
large for $N<d$, and at $N=d$ both summands must equal $\cos(2\pi/d)$.
When $n=2d$ in LT with $d$ odd, send $\zeta$ to the primitive root
$e^{\ii\pi(1-1/d)}$. On the even $N$-grid each remaining cosine is
at least $-\cos(2\pi/N)$. The target $-2\cos(\pi/d)$ is too small
for $N<2d$, and at $N=2d$ both must attain the minimum.

Pulling back gives the unique cosine pair $(E_*/2,E_*/2)$ at order
$n$. It yields the two constant configurations with angles
$\pm\arg\zeta$. On a reflected odd chart, adding $\pi$ to both
angles recovers the ordinary grid equation used here. Simultaneous
reflection equates the two chart integrals.

In the positive LT quadrant and the upper HT semicircle put
$\alpha=\beta=\theta$ in \eqref{iso:new:radial:L}. The two charts
give $2L_*$ in the $k$th coefficient derivative. The tail has bounded
derivatives through $k$ by Theorem~\ref{iso:new:budget}. The
argument of Theorem~\ref{iso:new:radial:full} gives the same radial
expansions with $\mathcal A_\eta=2\zeta_*^kp_\eta(\zeta_*)L_*$.
Evenness in LT, conjugation and HT inversion give the other points.
All these physical factors are analytic units.

For the stated prime powers the missing exponents are the multiples
of two or three, so the largest gap between primitive roots is
$4\pi/n$. Choosing the least permitted power strictly greater than
$4\pi/\ell$ gives the explicit bounds. For the asymptotic claim
take the least $n>4\pi/\ell$ in $n\equiv1\bmod6$ (HT) or
$n\equiv2\bmod12$ (LT), large enough that $n\ge 13$.
In HT every point of the $n$-grid except $1$ has an admissible odd
order; in LT every point except $\pm1$ has admissible order, and
its first even particle number divides $n$. The largest gap on
either retained grid is $4\pi/n$, and $n=4\pi/\ell+O(1)$.
A real $C^k$ extension would bound the $k$th derivative of the
radial restriction. Inversion and conjugation preserve angular
length.
\end{proof}

At a primitive $2^m$th root $e^{\ii\alpha}$ in the first quadrant,
$m\ge3$, put $E=s+s^{-1}$, $E_*=2\cos\alpha$, $k=2^{2m-1}-1$ and
$F(E)=X_L(s)/(1-s^{-4})^{1/4}$ on $|s|>1$. Then
\[
 \lim_{t\downarrow0}\sqrt t\,\partial_E^k F(E_*+\ii t)\ne0
\]
exists. Indeed the two reflected charts just found have the same
amplitude. In the local integral of Lemma~\ref{lem:local-period} one
can use $E-E_*$ as parameter: its derivative with respect to the
outward radial parameter at zero is $2\ii\sin\alpha$, so the vertical
ray has the same tangent as the radial approach. The local holomorphic
factor multiplying the half-integral power changes by an analytic unit.
Theorem~\ref{iso:new:budget} bounds the $k$th derivative of the higher
tail throughout the adjacent collar.

For a fixed odd prime $p\ge5$, every $2p$th root except $\pm1$ has
order $p$ or $2p$. Thus a low-temperature arc of length $>2\pi/p$
has an obstruction at derivative order $2p^2-1$. The unequal-angle
family alone gives that order on every arc of length $>2\sqrt{10}\pi/p$
when $p\ge11$. To see this put $h=2\pi/p$, $m=\lfloor(p-1)/4\rfloor$.
Adjacent grid angles $jh,(j+1)h$, $1\le j<m$, give
$\theta_j=\arccos[(\cos(jh)+\cos((j+1)h))/2]$. They increase and
lie between their grid angles. Concavity of cosine gives
$\theta_j\ge(j+1/2)h$, so consecutive gaps are at most $3h/2$.
Convexity of $\cos\sqrt u$ for $0\le u\le\pi^2$ gives
$\theta_1\le\sqrt{5/2}\,h$, and
$\pi/2-\theta_{m-1}\le5h/4$. Reflection in the axes bounds the
largest circular gap by $\sqrt{10}\,h$.

\section{The negative odd edge: boundary jets and their growth}
\label{neg:new:section}

At a negative edge, the odd momentum constraint prevents an endpoint
collision at every fixed particle index.  The complete sum is smooth on
the boundary, yet at the phases specified below its formal Taylor series
has radius of convergence zero. On the real angular plane, the two
smallest odd windings determine its coefficient growth. Throughout
this section $t$ denotes the edge coordinate $\tau$ in
Theorem~\ref{thm:unified}(vii).
Lemma~\ref{neg:new:transport} relates the indicator in the local edge
coordinate to its value in a physical transverse coordinate.

Fix $\phi=2\pi r/p\in(0,\pi)$, where $p\ge3$ is odd and
$(r,p)=1$, and write $b=\cos\phi$, $s_\phi=\sin\phi$.
Let $B(t,\eta)$ be holomorphic and bounded on a common disk times
a parameter polydisk $\mathcal P^+$, with
$\Im B(0,\eta)\le-\beta<0$.  All uniform assertions below are on
$\mathcal P\Subset\mathcal P^+$.  Shrinking the neighborhoods gives a common
logarithm of $B$, common scalar derivative bounds, and $\Im B<0$.
For fixed $\eta\in\mathcal P$, write $B(t)=B(t,\eta)$ and set
\[
 g(t,\delta)=\operatorname{arcosh}
       \{b+B(t)(t-1+\cos\delta)\},\quad g(0,0)=\ii\phi,
 \qquad \gamma(t,q)=g(t,q-\pi).
\]
The branch is analytic near all real angles. The transfer roots $e^{-g(0,\delta)}$
lie in a compact subset of the closed lower unit half disk, strictly
away from the real axis.  The physical side is
\[
             \mathcal U^- =\{|t|<r_0:\Im(B(t)t)>0\};
\]
it contains the inward ray $t=-\epsilon$.
For odd $n$ put
\[
 H_n(t,q)=\frac{B(t)^{n^2/2}}{n!(2\pi)^{n-1}}
 \prod_i\frac1{\sinh\gamma_i}
 \prod_{i<l}\frac{\sin^2((q_i-q_l)/2)}
                  {\sinh^2((\gamma_i+\gamma_l)/2)},\quad
 \Gamma_n=\sum_i\gamma_i,
\]
and $C_n(t)=\int_{\sum q_i=0\ (2\pi)}H_n\coth(\Gamma_n/2)\,\dd q$,
where $\dd q=\dd q_1\cdots\dd q_{n-1}$. Set $F=\sum_{n\ {\rm odd}}C_n$.

\begin{theorem}[Exact full-index indicator]\label{neg:new:indicator}
The complete sum is holomorphic on $\mathcal U^-$.  Every inward jet
$a_j=\lim_{\epsilon\downarrow0}F^{(j)}(-\epsilon)/j!$ exists, and
\begin{align}
 |a_j|&\le\exp\left\{j\left(\log j-\log\log j+
                         \log\frac4{\pi^2}-1\right)+o(j)\right\},
                                                        \label{neg:new:allj}\\
 \limsup_{j\to\infty}\frac{\log j}{j}|a_j|^{1/j}
      &=\frac4{e\pi^2}.                                \label{neg:new:type}
\end{align}
For the fixed phase, the upper error is uniform on compact admissible
holomorphic parameter families. For each family member and every
sufficiently large odd multiple $n$ of $p$, there is an integer $j_n$
such that
\[
 j_n=n^2\{\log n+\tfrac12\log\log n-c_*+o(1)\},\qquad
 c_*=-\tfrac34-\tfrac32\log2+\log\pi+
                       \tfrac12\log|B(0)|-\log s_\phi,
\]
and $|a_{j_n}|^{1/j_n}\sim4j_n/(e\pi^2\log j_n)$.
Thus no holomorphic extension through $t=0$ is possible.
\end{theorem}

Gaussian decay in the particle number gives every finite boundary jet.
For each resonant particle number, subtracting the pole leaves an
analytic remainder. A quadratic beta period gives the contribution
of the two smallest odd windings. Summation over the particle number
gives the upper bound, and along the sequence $j_n$ a single particle
number dominates and gives the matching lower bound.

\subsection{Boundary jets and the pole decomposition}

Lemma~\ref{edge:new:gaussian},
after the nonramified change
$h=b+B(t)t$, gives on a fixed complex $t$ disk
\begin{equation}
          |H_n(t,q)|\le E_n,\qquad
 E_n=\exp\{-\kappa n^2+Cn\log(n+1)+C\}.                  \label{neg:new:gauss}
\end{equation}
The constants include the torus volumes.  The momentum constraint now
gives a lower bound on the distance from the forbidden endpoint collision.
For shortest lifts
$q_i=\pi+\delta_i$, put $S=\sum_i\delta_i^2$.
Their sum is an odd multiple of $\pi$, hence $S\ge\pi^2/n$.
On the inward ray,
$\Re g(-\epsilon,\delta)\ge c(\epsilon+\delta^2)$.
A disk of radius $c(\epsilon+n^{-2})$ around $-\epsilon$ preserves
$\Re\Gamma_n\ge c(n\epsilon+n^{-1})$, since every parameter
derivative of an individual $g$ is bounded.  Therefore
\begin{equation}
 |[t^j]C_n|\le(Cn^2)^j CnE_n,\quad
 a_j=\sum_{n\ {\rm odd}}[t^j]C_n,\quad
 \sum_n|[t^j]C_n|<\infty.                              \label{neg:new:jets}
\end{equation}
Here $[t^j]C_n$ is the Taylor coefficient of the individual
holomorphic term at the edge. The $n$-dependent disks above justify
the sum of these coefficients for each fixed $j$.
Cauchy's formula on the parameter collar gives the same estimates for
each fixed parameter derivative.  Normal convergence and all the stated
boundary limits follow.

To isolate the terms responsible for large derivative order, choose a
small fixed $\sigma>0$.  On $S\ge\sigma$ the kernel is
analytic and bounded for $|t|<c/n$.  On $S<\sigma$, write
$D_n=\Gamma_n-n\ii\phi$; then $|D_n|\le C(S+n|t|)$.
If $p\nmid n$, the phase distance to $2\pi\mathbb Z$ is at least
$2\pi/p$, so this part is also analytic on the same kind of disk.
For odd multiples of $p$,
$\coth(\Gamma_n/2)=2/D_n+\mathcal R(D_n)$, with $\mathcal R$
holomorphic on a fixed disk.  The derivative
$n^{-1}\partial_tD_n$ lies in a fixed small disk about
$B(0)/(\ii s_\phi)\ne0$.  Rouch\'e's theorem on
$|t|=CS/n$ gives a unique zero $T_n(q)$, and
\[
 D_n=n(t-T_n)A_n(t,q),\qquad |T_n|\le CS/n,
\]
where $A_n$ is holomorphic on a fixed smaller disk and bounded away
from zero.  The latter facts also follow directly by integrating
$n^{-1}\partial_tD_n$ along the segment from $T_n$ to $t$.
For $J_n=H_n/A_n$ use the exact division
\[
 \frac{2H_n}{D_n}=
 \frac{2J_n(T_n,q)}{n(t-T_n)}+
 \frac2n\frac{J_n(t,q)-J_n(T_n,q)}{t-T_n}.
\]
The second term has a removable singularity at the pole and, together
with the $\mathcal R$ term, forms the analytic remainder.
Extracting coefficients, initially on $|t|<c/n^2$, gives
\begin{gather}
 a_j=\sum_{\substack{n\ {\rm odd}\\p\mid n}} A_{n,j}+R_j,\qquad
 A_{n,j}=-\int_{S<\sigma}
      \frac{2H_n(T_n,q)}{\partial_t\Gamma_n(T_n,q)}
                          T_n(q)^{-j-1}\,\dd q,\label{neg:new:poles}\\
 |R_j|\le \sum_n(Cn)^jE_n
       \le\exp\{\tfrac12j\log(j+2)+Cj+C\}.             \label{neg:new:remainder}
\end{gather}
For the last inequality, absorb $Cn\log(n+1)$ into half the Gaussian
exponent and maximize $j\log n-\kappa n^2/2$.

With $\xi_i=1-\cos\delta_i$ and $\bar\xi=n^{-1}\sum_i\xi_i$, the
linear terms in the averaged dispersion at $t=\bar\xi$ cancel.
The nonzero averaged derivative therefore gives
\begin{equation}
 \begin{gathered}
 T_n=\bar\xi+O\left(n^{-1}\sum_i(\xi_i-\bar\xi)^2\right),\\
 |T_n|\ge(1-C\sigma)S/(2n),\qquad
 |\partial_t\Gamma_n(T_n,q)|\ge cn.
 \end{gathered}                                      \label{neg:new:polelower}
\end{equation}
Dividing all sines by their linear terms and all hyperbolic sines by
$\ii s_\phi$ gives
\begin{equation}
 |H_n(T_n,q)|\le H_{0,n}\Delta(\delta)^2 e^{CnS},\qquad
 H_{0,n}=\frac{|B(0)|^{n^2/2}}
 {n!(2\pi)^{n-1}2^{n(n-1)}s_\phi^{n^2}}.              \label{neg:new:roughdensity}
\end{equation}
The logarithmic error from the anisotropy factor is bounded by
$Cn^2|T_n|\le CnS$. The one-body phase deviations sum in absolute
value to $O(S)$, giving a bound of $CnS$ for the logarithmic error
from the pair denominators. The sine divided differences are controlled
by $\sum_{i<l}|\delta_i-\delta_l|^2\le nS$.

Write a winding as
$\delta_i=w\pi/n+v_i$, $\sum_i v_i=0$, $w$ odd.  Put
\[
 Q=\sum_i v_i^2,\quad m_{n,w}=\pi^2w^2/n,\quad
 \alpha_n=(n^2-1)/2,\quad P_n=\prod_{r=1}^{n-1}\frac{r!}{2^r},
 \quad G_n=\frac{n!\pi^{(n-1)/2}P_n}{\sqrt n}.
\]
Thus $S=Q+m_{n,w}$ and $\Delta(\delta)=\Delta(v)$.
There are at most $C\sqrt n$ allowed windings. By \eqref{eq:shared-beta}
with $\ell=N=n$, their moments use $G_n$ from \eqref{eq:shared-G}.
Equations~\eqref{neg:new:polelower} and~\eqref{neg:new:roughdensity}
give, for $j+1>\alpha_n$, the winding bound
\begin{equation}
 |A_{n,j,w}|\le\frac{CH_{0,n}}n e^{Cn\sigma}
 \left(\frac{2n}{1-C\sigma}\right)^{j+1}
 G_nm_{n,w}^{\alpha_n-j-1}
       \frac{\Gamma(j+1-\alpha_n)}{\Gamma(j+1)}.        \label{neg:new:winding}
\end{equation}
For a bound valid at every jet index $j$, use $|T_n|\ge c/n^2$
and $\int_{Q<\sigma}\Delta^2\dd v
=G_n\sigma^{\alpha_n}/\Gamma(\alpha_n+1)$.  Since
$\log P_n=\tfrac12n^2\log n+O(n^2)$ and
$\log\Gamma(\alpha_n+1)=n^2\log n+O(n^2)$, this proves
\begin{equation}
 \log\sum_w|A_{n,j,w}|
 \le2j\log n-\tfrac12n^2\log n+Cn^2+Cj+C.             \label{neg:new:allindex}
\end{equation}

For large $j$ put $J=j/\log j$ and $T_j=j(\log J+\log(4/\pi^2)-1)$.
For some fixed $0<\eta_*<1/8$ it suffices to keep
\begin{equation}
                   \eta_* J\le n^2\le J/\eta_*.       \label{neg:new:middle}
\end{equation}
Indeed \eqref{neg:new:allindex}, with $x=n^2$, gives
$j\log x-\tfrac14x\log x+Cx+Cj+C$. For $x\le\eta_*J$ this is
at most $j\log J+j(\log\eta_*+C+o(1))$.
For $x\ge J/\eta_*$ we have $\tfrac14\log x-C\ge(\log j)/5$.
Writing $v=x/J$ gives the upper bound $j\log J+j(\log v-v/5+C)+C$.
Choose $\eta_*$ small enough that the first bound is below $T_j-2j$
and the second below $T_j-2j-x\log j/10$ for all large $j$.
The small indices number $O(\sqrt J)$, and the final Gaussian is
summable over the large ones. Together with \eqref{neg:new:remainder}
this bounds the total discarded moduli and $|R_j|$ by $O(e^{T_j-j})$.

\subsection{Residue concentration in the first two windings}

Reflection $\delta\mapsto-\delta$ interchanges the windings $w=\pm1$.
Their pole integrals agree, so it suffices to evaluate the one for $w=1$.

\paragraph{The pole and its residue.}
Set $d_n=\pi/n$, $\tau_n=1-\cos d_n$, $B_n=B(\tau_n)$, and
\[
 h_n=\cos d_n-\frac{B_nb}{s_\phi^2}\sin^2d_n=1+O(n^{-2}),
 \qquad a_n^*=\Re h_n>0.
\]
In the first winding $\delta_i=d_n+v_i$, $\sum_i v_i=0$.
The real integration region is the ball $Q<r_n^2=\sigma-\pi^2/n$.
Its radius is bounded above and away from zero for large $n$.
The implicit equation $\sum_i g(T_n,d_n+v_i)=n\ii\phi$ has a
holomorphic solution on a common small complex Euclidean ball,
with $|T_n-\tau_n|\le C\sum_i|v_i|^2/n$.
This follows by canceling the constant and linear shape terms and
applying Rouch\'e's theorem to the averaged equation; its parameter
derivative stays bounded away from zero on the fixed neighborhoods.
At the center
\[
 g_t(\tau_n,d_n)=-\ii B_n/s_\phi,\qquad
 \tfrac12g_{\delta\delta}(\tau_n,d_n)
                         =\ii B_nh_n/(2s_\phi).
\]
The $B'$ term in $W_t$, the derivative of the dispersion in braces,
vanishes exactly there.
Since $g$ is even in its angular argument,
$g_{\delta\delta\delta}(t,d_n)=O(d_n)$.  Taylor expansion on the
real ball, retaining the fourth-power sum, gives
\begin{align}
 T_n(v)&=T_{0,n}(v)+\mathcal R_n(v),
 &T_{0,n}&=\tau_n+h_nQ/(2n),\nonumber\\
 |\mathcal R_n(v)|&\le\frac Cn
 \left(n^{-1}\sum_i|v_i|^3+\sum_i v_i^4+Q^2/n\right).
                                                        \label{neg:new:realerror}
\end{align}
For the remainder estimate, write $u=T_n-\tau_n$.  The pole equation
then has remainder $O(n|u|^2+|u|Q+d_n\sum|v_i|^3+\sum v_i^4)$;
divide by its derivative of size $n$ and use
$|u|\le CQ/n$. Since $\sum|v_i|^3\le\sqrt{\sigma}\,Q$ and
$\sum v_i^4\le\sigma Q$, decreasing $\sigma$ gives
$|\mathcal R_n|\le a_n^*Q/(8n)$ for large $n$. Thus
\[
 D_{+,n}=\tau_n+a_n^*Q/(2n)\le|T_{0,n}|,\qquad
 |\mathcal R_n|\le D_{+,n}/4,\qquad
 \Re T_n\ge\tau_n+a_n^*Q/(4n).
\]

For the complete residue density
$\varrho_n=-2H_n(T_n,q)/\partial_t\Gamma_n(T_n,q)$ one has
\begin{align}
 \varrho_n(v)&=\varrho_{n,0}\Delta(v)^2e^{\ell_n(v)},
 &|\ell_n(v)|&\le CnQ,\label{neg:new:density}\\
 \varrho_{n,0}&=-\frac{2B_n^{\alpha_n-1/2}}
 {n\,n!(2\pi)^{n-1}2^{n(n-1)}s_\phi^{2\alpha_n}}.      \label{neg:new:residue}
\end{align}
The identities $\ii^{-n^2}=-\ii$ and
$\partial_t\Gamma_n(\tau_n,d_n)=-\ii nB_n/s_\phi$ fix the
phase and sign.  Reflection gives a second contribution with the same
value.  To estimate the logarithmic error, put $e_i=g(T_n,d_n+v_i)-\ii\phi$.
Then $\sum_i e_i=0$ exactly and
\[
 \sum_i|e_i|^2\le C(Q^2/n+Q/n^2+\sum_i v_i^4).
\]
The logarithm of $\sinh(\ii\phi+e)/(\ii s_\phi)$ is analytic.
Its linear terms cancel both in the one-body sum and in the paired
sum with arguments $(e_i+e_l)/2$. The remaining logarithmic error
is bounded by $Cn\sum|e_i|^2\le CnQ$. For the sine divided differences
the corresponding bound is $CnQ$, and for the anisotropy factor it is
$Cn^2|T_n-\tau_n|\le CnQ$. The normalized derivative sum is $1+O(Q/n)$.
These bounds
prove~\eqref{neg:new:density}, including at all coincidences.

\paragraph{The fourth angular moment.}
On the unit sphere of $H_0=\{\sum_i x_i=0\}$ let $\theta$ have density
proportional to $\Delta_n(\theta)^2$ in surface measure. Put
$P_4(x)=\sum_i x_i^4$ and $W_n=nP_4(\theta)$. For $v=\sqrt Q\,\theta$,
Cauchy--Schwarz gives
\begin{equation}
 \sum_i v_i^4=Q^2W_n/n,\qquad
 \sum_i|v_i|^3\le Q^{3/2}\sqrt{W_n/n}.                 \label{neg:new:powers}
\end{equation}
We bound $W_n$ by integration on the unit ball of $H_0$.
The tangent field $V_i=x_i^3-n^{-1}\sum_j x_j^3$ satisfies
$x\cdot V=P_4(x)$ and
\[
 \operatorname{div}_{H_0}(\Delta_n^2 V)=(2n-3/n)Q\,\Delta_n^2,
 \qquad Q=\sum_i x_i^2 .
\]
Indeed the tangential divergence of $V$ is $3(1-1/n)Q$.
Differentiating the squared differences gives the additional factor
$2\sum_{i<j}(x_i^2+x_ix_j+x_j^2)=(2n-3)Q$.
On the sphere, $\nabla P_4\cdot V\le4\sum_i x_i^6\le4P_4$.
Apply the divergence theorem to $P_4^{k-1}\Delta_n^2 V$, $k\ge1$.
Its flux density on the sphere is $P_4^k\Delta_n^2$. Radial
integration of the homogeneous divergence divides its sphere integral
by $n^2-1+4k-2$. Hence
\[
 \mathbb E W_n^k\le
 \frac{n(2n-3/n+4(k-1))}{n^2-1+4k-2}\,\mathbb E W_n^{k-1}
 \le(2+4(k-1)/n)\,\mathbb E W_n^{k-1}.
\]
The binomial series gives $\mathbb E e^{\delta W_n}\le
(1-4\delta/n)^{-n/2}$ for $n>4\delta$. This is uniformly bounded
for $n\ge8\delta$, and $W_n\le n$ covers the finitely many smaller
dimensions. Applying the same estimate with $\delta+1$ gives
\begin{equation}\label{neg:new:moment}
 \sup_{n\ge3}\mathbb E[(1+W_n)e^{\delta W_n}]<\infty
 \qquad(\delta>0).
\end{equation}
Homogeneity gives this angular law for every positive radial weight
in the beta period.

\paragraph{Comparison with the beta period.}
For $\Re h>0$ and $j+1>\alpha_n$, \eqref{eq:shared-beta} gives
\begin{equation}
 J_h=\int\frac{\Delta(v)^2\dd v}{(\tau_n+hQ/(2n))^{j+1}}
 =G_n(2n)^{\alpha_n}h^{-\alpha_n}\tau_n^{\alpha_n-j-1}
                     \frac{\Gamma(j+1-\alpha_n)}{\Gamma(j+1)}.
                                                        \label{neg:new:complexbeta}
\end{equation}
Its branch is fixed by $\Re h>0$.  Under the stated convergence
condition, all real-variable integrations are absolutely convergent.  Since $\Im h_n=O(n^{-2})$,
\begin{equation}
                   J_{a_n^*}/|J_{h_n}|=e^{O(n^{-2})}.   \label{neg:new:majorantcost}
\end{equation}

The first winding equals $\varrho_{n,0}J_{h_n}(1+O(n/\sqrt j))$
on \eqref{neg:new:middle}. More generally the relative error is
$O(n/\sqrt j+e^{-c n^2})$ as $n\to\infty$ and $n^2/j\to0$.

To prove this relative estimate, normalize the positive measure for
$J_{a_n^*}$ and put
$x=a_n^*Q/(2n\tau_n)$ and $\epsilon=n^2/j$.
Its radial law is proportional to $x^{\alpha_n-1}(1+x)^{-j-1}$,
independent of the angular variable $\theta$, and $Q\asymp x/n$.
Equations~\eqref{neg:new:realerror} and~\eqref{neg:new:powers} give
\[
 (j+1)|\mathcal R_n|/D_{+,n}\le
 \frac{Cj}{n^2}\frac{x^{3/2}\sqrt{W_n}+x^2(W_n+1)}{1+x}.
\]
Factor $\varrho_{n,0}\Delta(v)^2$ out of the exact integrand for
$w=1$ and its quadratic approximation. By~\eqref{neg:new:density} and
$|\mathcal R_n/T_{0,n}|\le1/4$, the remaining factors differ in
absolute value by at most
\begin{equation}
 D_{+,n}^{-j-1}\left[
 \exp\{CnQ+2(j+1)|\mathcal R_n|/D_{+,n}\}-1\right].          \label{neg:new:integranderror}
\end{equation}
After integration the common residue cancels from the relative error,
and the change from the positive beta period to $|J_{h_n}|$ costs
the factor in \eqref{neg:new:majorantcost}.
For $x\le32\epsilon$ the exponent is at most
$C\sqrt\epsilon(1+W_n)$. Since $e^v-1\le v e^v$ for $v\ge0$,
\eqref{neg:new:moment} bounds the angular mean of the bracket by
$O(\sqrt\epsilon)$ there. This central region lies inside $Q<r_n^2$
for all large $n$. On the rest of the ball,
$|T_n|\ge\Re T_n\ge\tau_n(1+x/2)$ and $|e^{\ell_n}|\le e^{Cn\sigma}$.
The quadratic model has the same lower denominator bound on its
entire tail. With $Z_\beta(\alpha,J')=\Gamma(\alpha)
\Gamma(J'+1-\alpha)/\Gamma(J'+1)$, substituting $u=x/2$ bounds both
tail integrals relative to the positive period by
\[
 C e^{Cn\sigma}\,2^{\alpha_n}(1+16\epsilon)^{-j/2}
       \frac{Z_\beta(\alpha_n,j/2)}{Z_\beta(\alpha_n,j)}
 \le C e^{Cn\sigma}\,8^{\alpha_n}(1+16\epsilon)^{-j/2}
 \le C e^{-c n^2}.
\]
The beta ratio is at most $4^{\alpha_n}$ by its product formula for
integer $\alpha_n\le j/4$. The last bound follows from
$j\log(1+16\epsilon)\sim16n^2$. Equation~\eqref{neg:new:majorantcost}
transfers the error to the complex period. The error is $O(n/\sqrt j)$
also throughout $\log j<n<\sqrt j/(\log j)^{1/4}$.

Reflection $\delta_i\mapsto-\delta_i$ preserves the dispersion,
squared sines and independent real measure.  Thus the first two
windings add with the same full complex phase:
\begin{gather}
 A_{n,j,1}+A_{n,j,-1}=M(n,j)(1+e_{n,j}),\qquad
 |e_{n,j}|\le Cn/\sqrt j,\label{neg:new:modelerror}\\
 M(n,j)=-\frac{2^{2-\alpha_n}B_n^{\alpha_n-1/2}P_n
                          n^{\alpha_n-3/2}}
        {\pi^{(n-1)/2}s_\phi^{2\alpha_n}}
 h_n^{-\alpha_n}\tau_n^{\alpha_n-j-1}
                \frac{\Gamma(j+1-\alpha_n)}{\Gamma(j+1)}.
                                                        \label{neg:new:model}
\end{gather}
\subsection{The upper envelope and a matching subsequence}
\label{sec:jet-subsequence}

Put $L(n,j)=\log|M(n,j)|$.  Uniform Stirling estimates give
\begin{equation}
 L(n,j)=2j\log n-\tfrac12n^2\log j+c_*n^2
       -j\log(\pi^2/2)+O(n\log n+n^4/j+j/n^2).         \label{neg:new:logmodel}
\end{equation}
Indeed
$\log P_n=\tfrac12n^2\log n-(\tfrac34+\tfrac12\log2)n^2
+O(n\log n)$,
$\log\tau_n=\log(\pi^2/2)-2\log n+O(n^{-2})$, and the gamma
logarithm is $-\alpha_n\log j+O(\alpha_n^2/j)$. Collecting the
terms of order $n^2$ gives $c_*$ in Theorem~\ref{neg:new:indicator}.

\paragraph{The all-index upper bound.}
For an arbitrary large $j$, each error, including $c_*n^2$, is
$o(j)$ throughout~\eqref{neg:new:middle}.  With $x=n^2$, the
leading expression is
$j\log x-\tfrac12x\log j-j\log(\pi^2/2)$.
It has its unique maximum at $x=2j/\log j$, with value $T_j$.
Summing over the $O(\sqrt J)$ indices in \eqref{neg:new:middle}
costs $e^{o(j)}$. Outside this range the pole terms and the analytic
remainder have already been bounded by $O(e^{T_j-j})$.

For the higher odd windings in \eqref{neg:new:middle}, replace in
$|M(n,j)|$ the factors $B_n,h_n,\tau_n$ by
$B(0),1,\pi^2/(2n^2)$ and call the resulting positive expression
$M_0(n,j)$.  Exact logarithms give
$|\log(M_0/|M|)|\le C(1+j/n^2)$.
Comparison with~\eqref{neg:new:winding} yields
\[
 |A_{n,j,w}|\le CM_0(n,j)e^{Cn\sigma}(1-C\sigma)^{-j-1}
                         |w|^{-2(j+1-\alpha_n)}.
\]
Choose $\sigma$ with $-\log(1-C\sigma)<\tfrac12\log3$.
Since $\alpha_n=o(j)$, $n=o(j)$ and $j/n^2=o(j)$ uniformly in
\eqref{neg:new:middle}, summing the at most $C\sqrt n$ windings with
$|w|\ge3$ gives $|M(n,j)|e^{-cj}$.
Together with~\eqref{neg:new:modelerror}, these estimates cover every
odd particle index.  The triangle inequality gives
\eqref{neg:new:allj}.

\paragraph{A dominant subsequence.}
Fix a large resonant index $n$ and put $d=2p$.
Extend $\alpha_x=(x^2-1)/2$ to real $x$. For $j\asymp n^2\log n$,
$m\asymp n$ resonant and $x=m+d/2$, the adjacent terms of
\eqref{neg:new:model} satisfy
\begin{equation}\label{neg:adjacent}
 \frac{L(m+d,j)-L(m,j)}{d}=G_j(x)+O_d(1),\qquad
 G_j(x)=\frac{2j}{x}-x\log(j-\alpha_x)+2c_*x .
\end{equation}
Indeed, the factorial increment divided by $d$ is the average of
$\log(r!)-r\log2$, $m\le r<m+d$. By Stirling's formula this is
$x\log x-(1+\log2)x+O_d(1)$, since the mean of $r+1/2$ is $x$.
The increment of $\log\Gamma(j+1-\alpha_m)$ is the negative sum
of $\alpha_{m+d}-\alpha_m=dx$ logarithms. Their arguments differ
from $j-\alpha_x$ by $O_d(x)$, giving
$-x\log(j-\alpha_x)+O_d(x^2/j)$ after division by $d$.
In the other factors it suffices to take the midpoint derivative of
\[
 (2j+1-m^2/2)\log m+
 \alpha_m\log\!\left(\frac{\pi^2|B(0)|}{4s_\phi^2}\right).
\]
The omitted linear term has bounded derivative. The corrections in
$\log|B_m/B(0)|$, $\log|h_m|$ and $\log(2m^2\tau_m/\pi^2)$
have convergent series in $m^{-2}$ with zero constant term.
These corrections and the midpoint derivative error cost
$O_d(1+j/m^3)$. Combining the terms gives \eqref{neg:adjacent}.

Choose $j_n$ as the nearest integer to a real solution of $G_j(n)=0$.
At $j=n^2(\log n+\tfrac12\log\log n-c_*\pm\epsilon)$ one has
$G_j(n)/n=\pm2\epsilon+o(1)$ for every fixed $\epsilon>0$.
On this interval $\partial_jG_j(n)=2/n-n/(j-\alpha_n)\asymp1/n$.
Hence the solution exists, is unique there, and rounding gives
$G_{j_n}(n)=O(1/n)$ and the asserted expansion for $j_n$.

At $j=j_n$, $n^2\log j_n/j_n\to2$. All the indices $m$ in
\eqref{neg:new:middle} therefore satisfy $m\asymp n$. Throughout
this interval
\[
 \partial_xG_j=-2j/x^2-\log(j-\alpha_x)+x^2/(j-\alpha_x)+2c_*
 \le-c\log j .
\]
Summing \eqref{neg:adjacent} from $n$ to $n+kd$ gives
$L(n+kd,j_n)-L(n,j_n)\le-c'd^2 k^2\log j_n+O_d(|k|)$.
Consequently all other models in the interval have total modulus
$o(|M(n,j_n)|)$. The relative errors tend uniformly to zero, and the
higher windings have the additional factor $e^{-cj_n}$. The discarded
indices and analytic remainder have total modulus $O(e^{T_{j_n}-j_n})$.

Finally~\eqref{neg:new:logmodel}
at the selected index gives $L(n,j_n)=T_{j_n}+o(j_n)$.  Thus
\[
                        a_{j_n}=M(n,j_n)(1+o(1)).
\]
This gives the lower bound in \eqref{neg:new:type}. The common
neighborhoods, bounds on $B$ and strict imaginary margin give
compact-family uniformity in the upper bound. At fixed $j$, Cauchy's
formula in the family parameter preserves this uniformity.

\subsection{Coordinate changes and physical normalization}

The indicator is invariant under multiplication by a nonvanishing
analytic prefactor. Under an analytic change of transverse coordinate
it scales by the absolute value of the derivative at the edge.

\begin{lemma}[Analytic coordinate and factor transport]
\label{neg:new:transport}
For $A(t)=\sum a_jt^j$ define
$\mathfrak t(A)=\limsup_j(\log j/j)|a_j|^{1/j}$.
If this is finite and positive, $g$ is analytic with $g(0)\ne0$,
and $\psi$ is analytic with $\psi(0)=0$, $\psi'(0)\ne0$, then
\[
 \mathfrak t(gA)=\mathfrak t(A),\qquad
 \mathfrak t(A\circ\psi)=|\psi'(0)|\mathfrak t(A).
\]
Adding an analytic series does not change $\mathfrak t$.
\end{lemma}
\begin{proof}
For $T>\mathfrak t(A)$, $|a_m|\le C(Tm/\log m)^m$ for $m\ge3$.
On $|t|\le r_j=(1-\varepsilon)\log j/(Tj)$, the truncation
$A_{\le j}$ is bounded independently of $j$, since $m/\log m$ is
increasing and each term is bounded by $C(1-\varepsilon)^m$.
The degree-$j$ coefficient of $gA$ equals that of $gA_{\le j}$.
Cauchy's estimate gives its modulus at most $C_\varepsilon r_j^{-j}$.
Letting $T$ decrease and $\varepsilon$ tend to zero proves the upper
inequality, also for input indicator zero.  Apply it to $1/g$ to
obtain equality.  For composition use
$r_j=(1-\varepsilon)\log j/(T|\psi'(0)|j)$;
$|\psi(t)|\le|\psi'(0)|r_j(1+O(r_j))$ keeps
$A_{\le j}\circ\psi$ uniformly bounded on that circle.  The same
argument proves the upper bound; composition with the analytic inverse
gives equality.  An analytic series has zero indicator, and addition
followed by subtraction proves the last assertion.
\end{proof}

For a nonramified physical transverse coordinate $z$ at a
selected negative endpoint $z_0$,
$t=(C+B-b)/B$ and $t'(z_0)=(C+B)'(z_0)/B(z_0)$.
Consequently the full reduced HT function, and its physical energy-field
normalization wherever that factor is nonzero, satisfy
\[
 \limsup_{j\to\infty}\frac{\log j}{j}
 \left|\frac{f^{(j)}(z_0)}{j!}\right|^{1/j}
 =\frac4{e\pi^2}\left|\frac{(C+B)'(z_0)}{B(z_0)}\right|.
\]
These are the continuous one-sided jets furnished by
Theorem~\ref{edge:new:budget}.  Their positive indicator makes the
formal Taylor series have radius zero. By Cauchy's estimate, a
holomorphic extension along any such curve would give a positive radius.

\section{Boundary jets and differential transcendence}
\label{da:new:section}

For each finite order, the preceding estimates give continuous jets
through that order on a nonempty open boundary arc. The local
obstructions are dense on this arc. The separant criterion uses both
properties to exclude algebraic differential relations for each
complete physical sum.

A differential identity of a germ with rational coefficients propagates
to its connected holomorphic physical domain: clear the finitely many
rational denominators and apply the identity theorem to the resulting
holomorphic polynomial in the complete jets.



\begin{theorem}[Separant criterion]\label{da:new:separant}
Let $f$ be holomorphic on a connected plane domain $D$. Suppose
that for every $r\ge0$ some nonempty relatively open real-analytic
boundary arc has a one-sided collar on which $f,\ldots,f^{(r)}$
extend continuously and finitely to the arc. If nonextendable
points from that collar are dense on each such arc, then $f$ is
differentially transcendental over $\C(z)$.
\end{theorem}
\begin{proof}
Clear denominators in a proposed relation $P(z,f,\ldots,f^{(r)})=0$.
Choose $r$ least and then minimize the positive degree in $Y_r$.
In characteristic zero the evaluated separant
$H=P_{Y_r}(z,f,\ldots,f^{(r)})$ is holomorphic and is not
identically zero, by the same minimality. It is continuous on
the chosen collar up to its boundary arc.

If $H$ is nonzero at a boundary jet, the holomorphic implicit
function theorem gives $f^{(r)}=G(z,f,\ldots,f^{(r-1)})$ near
that jet. For $r=0$ this continues $f$ directly. Otherwise set
$U=(f,\ldots,f^{(r-1)})$ and write the resulting first-order
system as $U'=V(z,U)$. On a smaller closed polydisk about the
boundary data, $V$ is bounded and Lipschitz in $U$. For initial
data in a concentric half-polydisk the local holomorphic solution
has a common positive radius. Indeed the integral form of the
system on a disk of radius $\delta$ is a contraction when
$\delta\sup\|D_U V\|<1$ and the displacement $\delta\sup|V|$
stays within the polydisk. Take an interior initial point so
close to the boundary point that its solution disk contains
that point. Uniqueness identifies the solution with the given
jet on the adjacent collar component, using a short inward segment
and the identity theorem. Thus a nonzero separant trace permits
continuation.

Consequently $H$ has zero trace on a dense subset of the arc,
and by continuity on the whole arc. Straighten a short analytic
subarc and use Schwarz reflection across the zero trace. The
identity theorem forces $H$ to vanish first there and then on
$D$, contrary to its choice.
\end{proof}

\section{Proof of the main theorem}\label{sec:assembly}

Proposition~\ref{geo:new:domains} gives (i), and
Corollary~\ref{geo:whole} identifies all other complementary components.
The angular representation \eqref{eq:global-seriesH}--\eqref{eq:global-series}
and its compact estimates prove (ii), with the physical normalization
fixed in Section~\ref{sec:physical}.

At a regular face use the chart $h=C-\sigma B$ of
Proposition~\ref{geo:new:faces}, after exchanging lattice axes
if necessary. A nonramified transverse curve has $B=B(h)$.
At each selected prime phase $b$, conjugate if needed to put the
physical side above the real $h$ axis. Proposition~\ref{prop:positive-edge}
gives the complete obstruction at $\sigma=+1$ in both parities.
At $\sigma=-1$ shifting every angle by $\pi$ gives the same result in
even parity. In odd parity $t=(h-b)/B(h)$ has derivative $1/B(b)$,
so Theorem~\ref{neg:new:indicator} and Lemma~\ref{neg:new:transport}
give the obstruction with its exact coordinate factor. The analytic
physical prefactors are units there. All three arguments admit the
specified compact parameter families with a strict imaginary margin.

These phases are dense on each regular face. A joint extension
would extend the constant-$B$ transversal. The density of regular
faces and Lemma~\ref{geo:new:collar} propagate this contradiction
to every finite boundary point and to any first exit of an unramified
continuation. This proves (iii).

For $\lambda\ne1$, \eqref{geo:new:pullbacks} gives the domain on the
temperature line. By Lemma~\ref{geo:new:slices}, its edge coordinate
satisfies $dh/dt\ne0$ at a dense set of boundary points with selected
phases. The transverse-curve conclusion of (iii) applies there, and
Lemma~\ref{geo:new:collar} gives nonextension at the remaining points.

For $\lambda=1$, Lemma~\ref{geo:new:cover} and
\eqref{eq:iso-physicalH}--\eqref{eq:iso-physical} give the disk,
exterior and covering normalizations. The nonzero full-sum periods
of Theorem~\ref{iso:new:radial:full} are dense on the circle, by
Lemma~\ref{iso:new:radial:arithmetic}, conjugation, HT inversion
and LT evenness. Their nonramified lifts are dense on the temperature
boundary. The collar argument includes its tips. An extension on
either physical side of any circle point would restrict to the
nearby selected points. The factor $t/(2J_x)$ is a unit at every
finite temperature boundary point. These observations give (iv)
and (v), including the magnetic energy-field responses.

For (vi), Theorem~\ref{edge:new:budget} bounds the tail on the closed
noncollinear collar under \eqref{eq:common-budget}. At isotropy,
Theorem~\ref{iso:new:budget} with $n=m+2$ gives the continuous
derivatives and their highest-order modulus under the same condition.
Given $r$, choose $m$ of the required parity with $2(r+1)<(m+2)^2-3$ and
remove the finite lower-index exception set from a noncritical arc.
Corollary~\ref{iso:new:jets} and Theorem~\ref{edge:new:budget}
give continuous jets on a smaller arc. At the negative odd endpoint
the momentum constraint gives all finite jets on the whole regular
arc. Analytic factors and nonramified changes of variable preserve
these finite jets by the chain rule. The expansions, finite parts,
and sharp radial moduli in (vii) are precisely those of
Theorem~\ref{iso:new:radial:full}. The remaining negative-edge
indicator, its upper estimate and its matching subsequence were
proved in Theorem~\ref{neg:new:indicator}.

Finally apply Theorem~\ref{da:new:separant} in each physical
variable itself. Parts (iv)--(vi) supply its two hypotheses for
$X_H,X_L$ in $s$, and for each $f_{\nu,\lambda}$ in $t$.
Multiplication by $t/(2J_x)$ preserves the finite jets on each
chosen arc and, being locally invertible as a factor there,
also preserves every local obstruction. The same criterion applies
to $\chi_{\nu,\lambda}$. A germ identity would first propagate
through the connected physical domain after its rational denominators
are cleared. This proves (viii) at every physical basepoint. On a
local isotropic $s$ chart the energy factor is a branch of
$\operatorname{arsinh}(s)/(2J_x)$. From a positive basepoint
continue that branch along a corridor to a noncritical open circle
arc avoiding $s=\pm i$ and $s=0$. It extends as a unit there.
For each prescribed order the finite jets exist after removal of
finitely many points on this arc. Theorem~\ref{da:new:separant}
therefore gives differential transcendence over $\C(s)$ for the
local susceptibility with its energy-field factor.\hfill\(\square\)

\begin{corollary}[The radial derivation]\label{da:radial}
Let $V=K_x\partial_{K_x}+K_y\partial_{K_y}$. Both $\mathcal X_H$ and
$\mathcal X_L$ are differentially transcendental over $\C(K_x,K_y)$
for the derivation $V$. More generally, the same conclusion holds
for a holomorphic function $F$ on a connected domain if
$F\circ\iota_\lambda$ has a differentially transcendental germ over
$\C(t)$ for some $\lambda>0$.
\end{corollary}
\begin{proof}
Fix such a $\lambda$ and clear denominators in a proposed polynomial
relation among $F,VF,\ldots,V^rF$. Divide its coefficients by their
largest common power of $K_y-\lambda K_x$. The resulting identity
holds throughout the domain by holomorphy. At least one coefficient
restricts to a nonzero polynomial on the line, since vanishing
identically there is equivalent to divisibility by $K_y-\lambda K_x$.
On $\iota_\lambda(t)$ the operator $V$ restricts to $t\,d/dt$.
Replacing its iterates by ordinary derivatives is triangular with
diagonal $1,t,\ldots,t^r$, invertible over $\C(t)$. The restricted
relation contradicts the differential transcendence of that germ.
For each physical susceptibility take $\lambda=1$ in
Theorem~\ref{thm:unified}(viii).
\end{proof}

\appendix
\section{Pair factors and periods at fixed total momentum}\label{app:forms}

\subsection{The pair product on a transfer curve}

In the dispersion relation $S(u+u^{-1})+T(z+z^{-1})=2A$,
put $H(v,w)=(v-w)/(1-vw)$. The complete pair in
\eqref{eq:split-pair} is a product of one factor in the angular
variables and one in the transfer roots. Iorgov and Lisovyy evaluate
the corresponding elliptic determinant for even $N$
\cite[(5.7)--(5.8), (6.4)]{IL}.

\begin{lemma}[A pointwise bound for the complete pair product]\label{lem:analytic-det}
Write $R_{ij}=H(u_i,u_j)$ and $P_{ij}=H(z_i,z_j)R_{ij}$
as in \eqref{eq:split-pair}, allowing limits from ordinary values.
Suppose $|u_i|\le M$, $|1-u_i u_j|\ge\delta>0$ for all $i,j$
(including $i=j$), and $|P_{ij}|\le M$. Then
\[
             \left|\prod_{i<j}P_{ij}\right|
                    \le C^N e^{-\kappa N^2}
\]
with constants depending only on $M,\delta$. If the pairs on
$\ell$ fixed labels have bounded holomorphic quotients
$P_{ij}/(z_i-z_j)^2$, the product after division by
$\prod_{i<j\le\ell}(z_i-z_j)^2$ is bounded by
$C_\ell^N e^{-\kappa N^2}$ on the same set, including coincidences.
Here $C_\ell$ may also depend on $\ell$ and the uniform quotient bound.
\end{lemma}
\begin{proof}
For a list $w=(w_1,\ldots,w_N)$ the ordinary Cauchy determinant gives
\[
 D(w):=\prod_{i<j}H(w_i,w_j)^2=\det[1+H(w_i,w_j)].
\]
Indeed $1+H(v,w)=(1+v)(1-w)/(1-vw)$. Clearing the denominators
in $\det[(1-v_iw_j)^{-1}]$ leaves the two Vandermonde factors by
alternation and degree. Comparison at $v_Nw_N=1$ fixes the constant
by induction. Hadamard's inequality gives $|D(w)|\le(C\sqrt N)^N$
where the entries are bounded.

On a fixed neighborhood of any permitted root tuple the stronger
bound $|D(w)|\le C^N e^{-cN^2}$ holds. Cover the bounded root set
by $s$ small disks, of radius chosen so that doubling any cluster
about its disk center preserves the reciprocal gaps. Some disk contains
$k\ge N/s$ roots. Dilate those roots about its
center by a common complex parameter $\lambda$. For $|\lambda|\le2$
the Hadamard bound applies, and $D$ has a zero of order at least
$k(k-1)$ at $\lambda=0$. The maximum principle at $\lambda=1$
gains $2^{-k(k-1)}$, which absorbs the factor $N^{N/2}$ after
decreasing $c$.

Choose a fixed $r>0$ so that the reciprocal gaps persist for
$|w_i-u_i|\le r$. If $|D(w)|\le B_N$ there, deleting up to $N$
factors $H(u_i,u_j)$ from $D(u)$ costs at most $C^N B_N$ in total
modulus, provided each pair is deleted at most twice. The remaining
product is holomorphic. Apply the submean inequality to its logarithmic
modulus on the product of circles $w_i=u_i+r e^{\ii\theta_i}$.
Each deleted numerator has mean log modulus at least $\log r$:
averaging in $\theta_i$ alone gives $\log\max(r,|u_i-w_j|)$.
Each denominator has bounded modulus. This proves the estimate,
including coincidences by continuity.

In each row of $[1+H(z_i,z_j)]$ choose an entry of largest modulus.
The diagonal entries are one. Each off-diagonal choice absorbs one
root factor, since
$|(1+H(z_i,z_j))H(u_i,u_j)|\le |H(u_i,u_j)|+|P_{ij}|\le C$.
No unordered pair is chosen more than twice. Hadamard's inequality
and the deletion estimate give
\[
 \left|\prod_{i<j}P_{ij}\right|^2=|D(z)D(u)|
      \le C^N N^{N/2} B_N .
\]
The bound extends to removable values by continuity. With the stronger
bound for $B_N$, decreasing the Gaussian exponent absorbs $N^{N/2}$.
For a fixed divided block apply this bound to the other $N-\ell$
labels, bound each cross pair by $M$ and each internal quotient by
its uniform bound.
\end{proof}

\subsection{Gaussian measure at fixed total angle}

For real $s$ use the coordinates $x_1,\ldots,x_{N-1}$ on the plane
$H_s=\{\sum_i x_i=s\}$ and write $\dd x'=\prod_{i<N}\dd x_i$.
Set $Q(x)=\sum_i x_i^2$ and
\[
 \Delta_\ell(x)=\prod_{1\le i<j\le\ell}(x_i-x_j),\qquad
 \alpha_{N,\ell}=\tfrac12\{N-1+\ell(\ell-1)\},\qquad
 P_\ell=\prod_{j=1}^{\ell-1}\frac{j!}{2^j}.
\]
Empty products equal one. The squared differences are invariant under
a common translation. This separates the total momentum from all the
relative coordinates.

\begin{lemma}[Periods at fixed total momentum]\label{lem:gaussian-period}
Let $N\ge2$, $0\le\ell\le N$, $\alpha=\alpha_{N,\ell}$ and
$G_{N,\ell}=\pi^{(N-1)/2}\ell!P_\ell/\sqrt N$.
For $t>0$,
\begin{equation}\label{eq:affine-G}
 \int_{H_s} e^{-tQ(x)}\Delta_\ell(x)^2\dd x'
 =G_{N,\ell}\,t^{-\alpha}e^{-ts^2/N}.
\end{equation}
Write $G_n=G_{n,n}$ and $a_n=\alpha_{n,n}=(n^2-1)/2$; in particular
\begin{equation}\label{eq:shared-G}
 G_n=\frac{n!\pi^{(n-1)/2}P_n}{\sqrt n}.
\end{equation}
For an integer $p>\alpha$, put $a_0=a+bs^2/N$. If $a_0,b$ belong
to a common open half-plane bounded by a line through zero, then
\begin{equation}\label{eq:shared-beta}
 \int_{H_s}\frac{\Delta_\ell(x)^2}{(a+bQ(x))^p}\dd x'
 =G_{N,\ell}\,b^{-\alpha}a_0^{\alpha-p}
             \frac{\Gamma(p-\alpha)}{\Gamma(p)} .
\end{equation}
Choose their arguments with difference in $(-\pi,\pi)$.
\end{lemma}
\begin{proof}
First compute the unrestricted integral at $t=1$. For
$d=\ell(\ell-1)/2$, write the Vandermonde as a determinant and replace
each variable in it by its partial derivative. Then
\[
 \Delta_\ell(\partial)e^{-Q}=(-2)^d\Delta_\ell(x)e^{-Q},
 \qquad \Delta_\ell(\partial)\Delta_\ell(x)
       =\ell!\prod_{j=0}^{\ell-1}j!.
\]
In the first determinant the $j$th derivative of the one-dimensional
Gaussian, divided by that Gaussian, has degree $j$ and leading
coefficient $(-2)^j$. In the second identity only matching
permutations survive differentiation. Integration by parts gives
$\int_{\R^N}e^{-Q}\Delta_\ell^2\,\dd x=\pi^{N/2}\ell!P_\ell$.
Write $x_i=\mu+v_i$ with $\sum v_i=0$. The Jacobian in
$(\mu,v_1,\ldots,v_{N-1})$ is $N$ and
$Q=N\mu^2+Q(v)$. Integration in $\mu$ therefore divides the full
integral by $\sqrt{N\pi}$ to give its value on $H_0$.
On $H_s$ the mean is $s/N$; scaling $v$ by $t^{-1/2}$
proves \eqref{eq:affine-G}.

On $H_0$, homogeneity gives the radial density
$G_{N,\ell}u^{\alpha-1}\dd u/\Gamma(\alpha)$ in $u=Q(v)$.
After translation to $H_s$ the denominator is $(a_0+bu)^p$.
The real beta integral gives \eqref{eq:shared-beta} for $a_0/b>0$
after factoring out $b^{-p}$. Both sides continue in $a_0/b$
on $\C\setminus(-\infty,0]$, since $|a_0/b+u|\ge c(1+u)$
on compact subsets there. This gives the stated branches.
\end{proof}

\subsection{Summation of the transfer powers}

Consider independent angles on the torus subject to
$\sum q_i=\vartheta$ modulo $2\pi$ and a candidate collision
at the common angle $q_0$. Put $q_i=q_0+x_i$ with
$-\pi\le x_i\le\pi$ and
$d_N=\dist(\vartheta-Nq_0,2\pi\Z)$. The lifts of the constraint
lie on the planes $H_{\vartheta-Nq_0+2\pi k}$.

\begin{proposition}[Derivatives of a transfer sum]\label{p:transfer}
Suppose a holomorphic transfer product $Z=\prod e^{-\gamma_i}$
in a rectangular collar $z=b+\ii y$, $b\in I$, $0<y<y_0$, satisfies
$\Re\gamma_i\ge c(y+x_i^2)$ throughout the collar.
Suppose the logarithms have uniformly
bounded first parameter derivatives on a common complex neighborhood of the closed collar.
For $N\ge n\ge2$, let the divided numerator $A_N(z,q)$ be
holomorphic in $z$ there, with $|A_N|\le C_n^N e^{-\kappa N^2}$.
Set $q_N=\vartheta-\sum_{i<N}q_i$ modulo $2\pi$. For $L\ge1$ put
\[
 U_{L,N}(z)=\frac{1}{N!}\int_{\mathbb T^{N-1}}
 A_N(z,q)\,Z(z,q)^L
 \prod_{i<k\le n}\sin^2\!\frac{q_i-q_k}{2}
 \prod_{i<N}\frac{\dd q_i}{2\pi},\qquad \mathbb T=\R/(2\pi\Z).
\]
For $\partial=\partial_z$ and fixed $j$ one has
\begin{equation}\label{eq:edge-transfer}
 |\partial^j U_{L,N}|
 \le\frac{C_{n,j}^N}{N!}e^{-\kappa N^2}
       L^{j-\alpha_{N,n}}e^{-c' L(Ny+d_N^2/N)} .
\end{equation}
The derivative of order $j$ in the sum over $L$ and $N\ge n$
extends continuously from inside when $2j<n^2-3$. At the greatest
such order the modulus is $O(\sqrt h)$ for even $n$ and
$O(h\log(e/h))$ for odd $n$ on a shorter closed collar.
\end{proposition}
\begin{proof}
On a parameter disk of radius $\epsilon/(1+L)$ the change of
$L\sum\gamma_i$ is $O(N)$. Replace the squared sines by a
constant times $\Delta_n^2$ and enlarge each lifted piece to its
whole plane. By \eqref{eq:affine-G} the shape integral contributes
$G_{N,n}(cL)^{-\alpha_{N,n}}$.
The sum of $\exp[-cL(\vartheta-Nq_0+2\pi k)^2/N]$ is bounded by
$C(1+\sqrt{N/L})e^{-cLd_N^2/(2N)}$, by factoring out half the
distance exponent and comparing the rest with a Gaussian integral.
Cauchy's estimate gives \eqref{eq:edge-transfer}.
The largest power is $j-(n^2-1)/2$, at $N=n$.
At the greatest admissible $j$ it is $-3/2$ or $-2$.
Integrate the $(j+1)$st derivative over a segment of length $h$
for $L\le1/h$ and use twice the $j$th derivative bound for
$L>1/h$. Summing gives the asserted moduli.
\end{proof}

\Needspace{14\baselineskip}
\subsection{The local Cauchy integral}

\begin{lemma}[A scalar integral for a fixed collision]\label{lem:local-period}
Let $n\ge2$. On $v_n=-\sum_{i<n}v_i$ use $\dd v=\prod_{i<n}\dd v_i$,
$Q=\sum v_i^2$ and $\Delta_n=\prod_{i<j}(v_i-v_j)$. Put
$G_n=\int e^{-Q}\Delta_n^2\dd v$, $a=(n^2-1)/2$,
$k=\lfloor a\rfloor$, $\nu=k+1-a$. Let $A,D$ be holomorphic near
$(0,0)$, with $A(0,0)=A_0$ and
\[
 D(z,v)=dz+hQ+O(|z|^2+|z|\,|v|+|v|^3),\qquad dh\ne0 .
\]
Let $\chi$ be a bounded cutoff equal to one near zero with sufficiently
small support. In a sector about a ray $z=\tau\epsilon$ on which
$d\tau$ and $h$ lie in a common open half-plane bounded by a line
through zero, write $c=h/d$ and $w=z/c$. Then
\[
 \int\frac{\chi(v)\Delta_n(v)^2 A(z,v)}{D(z,v)}\,\dd v
 =P(z)+\frac{B(z)}{c}
 \begin{cases}
   \displaystyle\frac{\pi}{\sin\pi a}\,w^{a-1},&a\notin\mathbb Z,\\
   (-1)^a w^{a-1}\log w,&a\in\mathbb Z ,
 \end{cases}
 \qquad B(0)=\frac{A_0G_n}{d\,\Gamma(a)} .
\]
Both $P$ and $B$ are holomorphic at zero. Take $-\pi<\arg w<\pi$.
The $k$th derivative with respect to $z$, evaluated at $z=\tau\epsilon$, is
\begin{equation}\label{eq:local-period}
 (-1)^k A_0d^k G_n h^{-a}(d\tau)^{-\nu}\Gamma(\nu)\epsilon^{-\nu}
  +\begin{cases}O(1),&\nu=1/2,\\ O(\log(e/\epsilon)),&\nu=1.\end{cases}
\end{equation}
Choose the arguments of $d\tau,h$ with difference $\arg(\tau/c)$.
Lower derivatives are bounded, except at order $k-1$ when $\nu=1$,
where the bound is $O(\log(e/\epsilon))$.
\end{lemma}
\begin{proof}
The implicit function theorem gives
$D=U(z,v)(z-\omega(v))$, where $U(0,0)=d$ and
$\omega(v)=-cQ+O(|v|^3)$. Changing the cutoff to the indicator of
a small ball about zero adds a holomorphic function, since the
denominator at $z=0$ is separated from zero on the intervening annulus.
In polar coordinates $v=r\eta$, $Q(\eta)=1$, the measure
$\Delta_n^2\dd v$ is $r^{2a-1}\dd r\,\dd m(\eta)$ with
$\int\dd m=2G_n/\Gamma(a)$, by integrating $e^{-r^2}$.

For each $\eta$ put $\xi=r\sqrt{-\omega(r\eta)/(cr^2)}$, choosing
the square root tending to one. This is holomorphically invertible
in $r$ on uniform small disks, with inverse $R(\xi,\eta)=\xi+O(\xi^2)$.
Its image of a short real radius lies in an arbitrarily narrow wedge
about the positive axis. On a sufficiently small sector about the
given ray, the two poles $\xi^2=-z/c$ stay outside this wedge.
Deform the image path to a common real segment plus an endpoint
segment away from zero. The latter contributes holomorphically at
zero because its denominator is separated from zero there.

On the real segment the numerator after $\xi^{2a-1}$ is removed is
$(A/U)(z,R\eta)(R/\xi)^{2a-1}\partial_\xi R$. Since
$R(-\xi,-\eta)=-R(\xi,\eta)$ and $\dd m(-\eta)=\dd m(\eta)$,
its spherical average is even in $\xi$ and holomorphic, by uniform
convergence on smaller disks. Substitution $u=\xi^2$ leaves
\[
 \int_0^b \frac{u^{a-1}B_1(z,u)}{z+cu}\,\dd u,\qquad
 B_1(0,0)=\frac{A_0G_n}{d\Gamma(a)},
\]
with $B_1$ holomorphic. Its difference from $B(z)=B_1(z,-z/c)$
is divisible by $u+z/c$, hence contributes another holomorphic
function. For $J_a(w)=\int_0^b u^{a-1}(w+u)^{-1}\dd u$ the
recurrence $J_{a+1}=b^a/a-wJ_a$ suffices. Modulo functions
holomorphic at zero, $J_1=-\log w$ and $J_{1/2}=\pi/\sqrt w$.
The latter follows by putting $u=t^2$ and subtracting the integral
over $t\ge\sqrt b$ from $\int_0^\infty2(w+t^2)^{-1}\dd t$.
Evaluation at real positive $w$ followed by continuation fixes the
branch. These formulas give the displayed decomposition.
Differentiation, using $\Gamma(1/2)=\sqrt\pi$ and
$(w^{k-1}\log w)^{(k)}=(k-1)!/w$, gives \eqref{eq:local-period}.
\end{proof}

\begin{thebibliography}{99}
\bibitem{ONGP}
W. P. Orrick, B. Nickel, A. J. Guttmann and J. H. H. Perk,
\emph{The susceptibility of the square lattice Ising model: New developments},
Journal of Statistical Physics 102(3--4) (2001), 795--841,
\href{https://doi.org/10.1023/A:1004850919647}{doi:10.1023/A:1004850919647},
\href{https://arxiv.org/abs/cond-mat/0103074}{arXiv:cond-mat/0103074}.

\bibitem{GE}
A. J. Guttmann and I. G. Enting,
\emph{Solvability of some statistical mechanical systems},
Physical Review Letters 76(3) (1996), 344--347,
\href{https://doi.org/10.1103/PhysRevLett.76.344}{doi:10.1103/PhysRevLett.76.344}.

\bibitem{Nickel}
B. Nickel, \emph{On the singularity structure of the 2D Ising model
susceptibility}, Journal of Physics A: Mathematical and General 32(21)
(1999), 3889--3906,
\href{https://doi.org/10.1088/0305-4470/32/21/303}{doi:10.1088/0305-4470/32/21/303}.

\bibitem{TW}
C. A. Tracy and H. Widom,
\emph{On the singularities in the susceptibility expansion for the
two-dimensional Ising model}, Journal of Statistical Physics 156(6)
(2014), 1125--1135,
\href{https://doi.org/10.1007/s10955-014-1061-4}{doi:10.1007/s10955-014-1061-4}.

\bibitem{MM}
B. M. McCoy and J.-M. Maillard,
\emph{The importance of the Ising model},
Progress of Theoretical Physics 127(5) (2012), 791--817,
\href{https://doi.org/10.1143/PTP.127.791}{doi:10.1143/PTP.127.791},
\href{https://arxiv.org/abs/1203.1456}{arXiv:1203.1456}.

\bibitem{GJMP}
A. J. Guttmann, I. Jensen, J.-M. Maillard and J. Pantone,
\emph{Is the full susceptibility of the square-lattice Ising model a
differentially algebraic function?},
Journal of Physics A: Mathematical and Theoretical 49(50) (2016), article 504002,
\href{https://doi.org/10.1088/1751-8113/49/50/504002}{doi:10.1088/1751-8113/49/50/504002},
\href{https://arxiv.org/abs/1607.04168}{arXiv:1607.04168}.

\bibitem{BL}
A. I. Bugrij and O. O. Lisovyy,
\emph{Correlation function of the two-dimensional Ising model on a finite
lattice: II}, Theoretical and Mathematical Physics 140(1)
(2004), 987--1000,
\href{https://doi.org/10.1023/B:TAMP.0000033035.90327.1f}{doi:10.1023/B:TAMP.0000033035.90327.1f},
\href{https://arxiv.org/abs/0708.3643}{arXiv:0708.3643}.



\bibitem{CGN}
F. Camia, C. Garban and C. M. Newman,
\emph{The Ising magnetization exponent on $\mathbb Z^2$ is $1/15$},
Probability Theory and Related Fields 160(1--2) (2014), 175--187,
\href{https://doi.org/10.1007/s00440-013-0526-8}{doi:10.1007/s00440-013-0526-8},
\href{https://arxiv.org/abs/1205.6612v4}{arXiv:1205.6612v4}.

\bibitem{BHM}
S. Boukraa, S. Hassani and J.-M. Maillard,
\emph{Holonomic functions of several complex variables and singularities
of anisotropic Ising n-fold integrals}, Journal of Physics A:
Mathematical and Theoretical 45(49) (2012), 494010,
\href{https://doi.org/10.1088/1751-8113/45/49/494010}{doi:10.1088/1751-8113/45/49/494010},
\href{https://arxiv.org/abs/1207.1784v2}{arXiv:1207.1784v2}.

\bibitem{IL}
N. Iorgov and O. Lisovyy,
\emph{Ising correlations and elliptic determinants},
Journal of Statistical Physics 143(1) (2011), 33--59,
\href{https://doi.org/10.1007/s10955-011-0154-6}{doi:10.1007/s10955-011-0154-6},
\href{https://arxiv.org/abs/1012.2856v2}{arXiv:1012.2856v2}.

\end{thebibliography}
\end{document}
